% ==========================================
% IMO 2026 Solution Notes - LaTeX - main.tex 
% Compiled by XeLaTeX
% Warning! Please do NOT direct copy and paste this code, because this is a older version demo. NOT for the new version.
% 請不要直接複製這段代碼並用於最新模板當中，因為這段代碼已經不適用於最新的模板包了。我後續會替換並展出一個新的Demo。
% ==========================================

\documentclass[11pt]{article}
\usepackage{jesse} 

\begin{document}

% ==========================================
% 首頁與裝飾設計
% ==========================================

\begin{titlepage}
    \begin{tikzpicture}[remember picture,overlay]
        \draw[GlacierBlue, line width=2pt] (current page.north west) ++(2cm,-1.5cm) -- ++(17.5cm,0);
        \draw[GlacierBlue!50, line width=0.6pt] (current page.north west) ++(2cm,-1.9cm) -- ++(17.5cm,0);
    \end{tikzpicture}
    
    \vspace*{3em}
    \begin{center}
        {\Huge \textbf{International Mathematical Olymiad}\\[1.5cm]}
        {\LARGE \textbf{Solution Notes}\\[1.5cm]}
        \textcolor{GlacierBlue}{\rule{0.5\textwidth}{1pt}} \\[1.75cm]
        {\Huge \calli{Jesse Zhang}} \\[10cm]
        {\LARGE \textbf{2026 Edition}}\\[1cm]
    \end{center}
    
    \vfill
    \begin{tikzpicture}[remember picture,overlay]
        \fill[GlacierBlue!20] (current page.south west) ++(2cm,2cm) rectangle ++(17.5cm,0.4cm);
        \draw[GlacierBlue!70!black, line width=1pt] (current page.south west) ++(2cm,2.2cm) -- ++(17.5cm,0);
    \end{tikzpicture}
\end{titlepage}

\tableofcontents
\newpage
\section{Solutions to Day I}
\subsection{International Mathematical Olympiad/Problem 1}
\begin{problem*}
    There are 2026 integers greater than 1 written on a blackboard, not necessarily different. In a move, Confucius chooses two integers $m > 1$ and $n > 1$ from different places on the blackboard and replaces these two integers with
\[
\gcd(m,n) \qquad \text{and} \qquad \frac{\text{lcm}(m,n)}{\gcd(m,n)}.
\]
He continues to make moves while it is possible to do so.

\begin{itemize}
    \item[(a)] Prove that, regardless of the choices of Confucius, after finitely many moves, exactly one integer $M$ on the blackboard is greater than 1.
    \item[(b)] Prove that the value of $M$ does not depend on the choices of Confucius.
\end{itemize}
\end{problem*}
\subsubsection{Solution 1}
Say we remove one element, call it $e$. Form the $45 \times 45$ matrix and call it our blackboard $\mathcal{B}$. At each step, choose two $m,n > 1$ such that they are replaced in $\mathcal{B}$ to\[\gcd(m,n)\quad\text{and}\quad\frac{\operatorname{lcm}(m,n)}{\gcd(m,n)}=\frac{mn}{\gcd(m,n)^2}.\]
Let $g = \gcd(m,n)$.
\begin{lemma}
    The 2025 elements in the matrix have the property that they will never degenerate to a case where all elements are 1.
\end{lemma}
\begin{proof}
    Suppose two entries $m,n > 1$ are inserted. Then, they are replaced by $g$ and $\frac{mn}{g^2}$. Assume that both replacements degenerate to 1. Then $g = 1$ from which $mn$ is the remaining. But $m,n > 1$ hence this is impossible. So at least one element after this process is greater than 1.\hfill$\square$
\end{proof}
\begin{lemma}
    Suppose $\mathcal{B}$ contains $k \geq 1$ elements greater than 1. If $e > 1$ is the removed element, then adding back $e$ and performing the same moveset as above degenerates the new blackboard to a case where exactly one element is greater than 1.
\end{lemma}
\begin{proof}
    Let $a,b > 1$ be two numbers in $k$. Consider their prime factorizations. For any prime $p$ let $v_p(a)$ and $v_p(b)$ be the power of $p$ in $a,b$ respectively. The moveset replaces the values with $\min(v_p(a),v_p(b))$ and $|v_p(a)-v_p(b)|$. However, this is just the Euclidean algorithm on the exponents of any given prime in $a,b$. The algorithm always ends with a nonzero value and one zero value for $v_p(a,b)$. Hence, one of the numbers is greater than one and the other is one itself. Therefore, exactly one number exists that is greater than 1.\hfill$\square$
\end{proof}
\paragraph{For Part (a):} \textcolor{BrightAqua}{\textbf{Lemma 1.1.2}} proves the result. $\square$

\medskip
\noindent
We now prove that $e$ is invariant and no matter what we choose the value of $e$ to be in the original 2026 matrix, the result always holds.
\begin{lemma}
$\gcd\bigl(v_p(a_1),v_p(a_2),\dots,v_p(a_n)\bigr)$ is invariant under any move.
\end{lemma}
\begin{proof}
    If Confucius chooses $m$ and $n$, then let $x=v_p(m), y=v_p(n)$. These moves replace $m$ and $n$ by $\gcd(m,n), \frac{\operatorname{lcm}(m,n)}{\gcd(m,n)}$. By Lemma 2, that is $\min(x,y)$ and $|x-y|$. In Lemma 2, we also showed this is the Euclidean Algorithm. Hence, $\gcd(x,y)=\gcd(\min(x,y),|x-y|)$, which is the same greatest common factor. Extending to $n$ such integers, the result follows.\hfill$\square$
\end{proof}
\paragraph{For Part (b):} By \textcolor{BrightAqua}{\textbf{Lemma 1.1.2}} exactly one number $M$ remains and all others are 1. Thus $$\gcd(M,1,1,\dots,1)=M.$$ In addition, the greatest common factor of all numbers on the blackboard is just $$\prod_{p} p^{v_{p}(M)} = M.$$ However, by \textcolor{BrightAqua}{\textbf{Lemma 1.1.3}} $p^{v_{p}(M)}$ for all $p$ is invariant hence $M$ is invariant. $\square$

\subsubsection{Solution 2}
\paragraph{For Part (a):} Let the $2026$ numbers be $x_1, x_2, \dots, x_{2026}$. For any two arbitrary numbers $x, y$ among these $2026$ numbers, let $d=\gcd(x,y)$. Then there exist coprime positive integers $a, b$ such that $x=da, y=db$. We then have:
$$\dfrac{\operatorname{lcm}(x,y)}{\gcd(x,y)}=\dfrac{\operatorname{lcm}(x,y) \cdot \gcd(x,y)}{\gcd(x,y)^2}=\dfrac{xy}{d^2}=\dfrac{x}{d} \cdot \dfrac{y}{d}=ab.$$Therefore, the two new numbers are $d$ and $ab$.

On the other hand, since 
\begin{align*}
g(x)+g(y)&=g(d \cdot a)+g(d \cdot b)\\&=g(a)+g(d)+g(b)+g(d)\\&=g(a)+g(b)+2g(d),
\end{align*}
and $$g(ab)+g(d)=g(d)+g(a)+g(b).$$Thus, when replacing $x, y$ with $d, ab$ in the sequence of $2026$ numbers, the sum $$S=g(x_1)+g(x_2)+...+g(x_{2026}).$$ decreases by exactly $g(d)$ at each step.

Let $C$ be the number of integers strictly greater than $1$ on the board. Each step replaces two numbers $x, y$ (both $>1$) with $d$ and $ab$, while the rest remain unchanged. The two removed numbers contribute exactly $2$ to $C$, whereas the two new numbers contribute at most $2$ to $C$, so $C$ does not increase after each step. Furthermore, the two new numbers cannot both be $1$: if $\operatorname{lcm}(x,y)/\gcd(x,y)=1$, then $\operatorname{lcm}(x,y)=\gcd(x,y)$, meaning $x=y$, in which case $d=\gcd(x,y)=x>1$, a contradiction. Therefore, at each step, $C$ decreases by either $0$ or $1$. Since a step can only be executed if there are at least two numbers $>1$, and after the step there is still at least one number $>1$, starting from $C=2026$, we always have $C\ge 1$ throughout the process.

Next, we set $t=S+C$. Since $t$ is a positive integer, we observe the following:
\begin{itemize}
\item If $d>1$, then $S$ decreases by at least $1$.
\item If $d=1$, then $x, y$ are coprime, the two new numbers are $1$ and $xy>1$, so $C$ decreases by exactly $1$ (while $S$ remains the same because $\Omega(d)=0$).
\end{itemize}

Thus, in each step, $t$ decreases by at least $1$, so the process must terminate after a finite number of steps. When it stops, no two numbers greater than $1$ can be chosen, meaning $C\le 1$; combining this with $C\ge 1$ yields $C=1$. So, exactly one number $M>1$ remains on the board.

\paragraph{For Part (b):}
For each prime number $p$, let $u=v_p(x),\ v=v_p(y)$. We have $$v_p(\gcd(x,y))=\min(u,v),\quad v_p(\operatorname{lcm}(x,y))=\max(u,v),$$
and$$v_p\left(\dfrac{\operatorname{lcm}(x,y)}{\gcd(x,y)}\right)=v_p(\operatorname{lcm}(x,y))-v_p(\gcd(x,y))=|u-v|.$$
Thus, when replacing $x, y$ with $\gcd(x,y), \frac{\operatorname{lcm}(x,y)}{\gcd(x,y)}$, the pair $(v_p(x), v_p(y))$ transforms into $$\left(v_p(\gcd(x,y)), v_p\left(\frac{\operatorname{lcm}(x,y)}{\gcd(x,y)}\right)\right).$$ Since $\gcd(u,v)=\gcd(\min(u,v),|u-v|)$, and the other numbers remain unchanged, we have an invariant: for each prime $p$, the greatest common divisor of the $v_p$'s across the entire board does not change through any steps.

In the final state, the board consists of $M$ and $2025$ ones. For each prime $p$, the exponents $v_p$ on the board are $v_p(M)$ along with $0$'s, so the greatest common divisor of the $v_p$'s across the board is simply $v_p(M)$. Therefore, $$v_p(M)= \gcd(v_p(x_1),v_p(x_2),...,v_p(x_{2026})),$$ meaning $$M=\prod_{p \mid M} p^{\gcd(v_p(x_1),v_p(x_2),...,v_p(x_{2026}))}.$$ The right-hand side depends only on the initial set of numbers and is independent of Confucius's choices. Furthermore, each $x_i>1$ has a prime divisor $p$, so the greatest common divisor of the $v_p$'s across the board is always greater than or equal to $1$ for at least one prime $p$, reaffirming that $M>1$. $\square$
\newpage
\subsubsection{Solution 3}
Let's fix initial $2026$ integers, and let those be $a_1,a_2,\dots , a_{2026}$. Let us construct a directed graph consisting of sets of $2026$ numbers and connect $A \to a$, where $A, a \in \mathbb Z_+^{2026}$ if and only if $A$ is achievable from $a$ by providing one $\gcd$-$\mathrm{lcm}$ operation. Let now switch to the exponent analysis, that is we change $a_i$ to $\nu_p(a_i)$ and observe our operation now is $(x,y) \to (|x-y|,\min(x,y))$. We will now use Newman's lemma to show the uniqueness of the terminal position for any $p$. Let’s show that, by carrying out any two steps, we can arrive at the same final result:
If $a \ge b \ge c$ in such for the following pair of operations:
\[
\begin{aligned}
(a,b,c\underbrace{...}) \to (a-b,b,c\underbrace{...}) \to (a-b,b-c,c\underbrace{...}) 
\\
(a,b,c\underbrace{...}) \to (b-c,b,c\underbrace{...}) \to (a-b,b-c,c\underbrace{...}) 
\end{aligned}
\]
If $a \ge b \ge c$ in such for the following pair of operations:
\[
\begin{aligned}
(a,b,c\underbrace{...}) \to (a-b,b,c\underbrace{...}) \to (a-b,a-c,c\underbrace{...}) 
\\
(a,b,c\underbrace{...}) \to (a-c,b,c\underbrace{...}) \to (a-b,b-c,c\underbrace{...}) 
\end{aligned}
\]
If $a \ge b, c \ge d$ in such for the following pair of operations:
\[
\begin{aligned}
(a,b,c,d\underbrace{...}) \to (a-b,b,c,d\underbrace{...}) \to (a-b,c-d,b,d\underbrace{...}) 
\\
(a,b,c,d\underbrace{...}) \to (a,b,c-d,d \underbrace{...}) \to (a-b,c-d,b,d\underbrace{...})
\end{aligned}
\]
Thus, if we remove the cycles from the graph where the position with $2025$ zeros comes to a cycle, from the unity of the final position we can take simply $p$ raised to the powers of the exponents of the sets, and it is clear that, since $$|x-y|+\min (x,y)=\max (x, y) > x+y, \quad \forall x,y>0,$$ for any prime $p$, eventually $2025$ valuations will be zero at the final position. Futhermore, returning to the initial sets, if two numbers have valuations greater than $0$ for any two primes, then the position is not terminal, as a move can still be made. Thus, eventually, for $2025$ numbers, the valuation modulo any prime is $0$, which completes part a. For the part b just note that from Euclidean algorithm $\gcd (x,y) = \gcd (|x-y|,\min (x,y))$. Hence, \(\gcd(\nu_p(a_1),\nu_p(a_2),\ldots,\nu_p(a_{2026}))\) is invarint and therefore:
$$M = \prod_{p\text{ prime}}^{}p^{\gcd(\nu_p(a_1),\nu_p(a_2),\ldots,\nu_p(a_{2026}))},$$which depends only on $a_1,a_2,\dots , a_{2026}$ as required. $\Box$

\newpage
\subsection{International Mathematical Olympiad/Problem 2}
\begin{problem*}
    Let $ABC$ be a triangle and let points $M$ and $N$ be the midpoints of sides $AB$ and $AC$, respectively. Let points $K$ and $L$ be chosen strictly inside triangles $BMC$ and $BNC$, respectively, such that $K$ lies strictly inside triangle $ABL$ and $L$ lies strictly inside triangle $AKC$. Suppose that
\[
\angle KBA = \angle ACL, \qquad \angle LBK = \angle LNC, \qquad \text{and} \qquad \angle LCK = \angle BMK.
\]
Let $O$ be the circumcentre of triangle $AKL$. Prove that $OM = ON$.
\end{problem*}
\subsubsection{Solution 1}
\begin{figure}[H]
    \centering
\begin{asy}
    size(14cm);
    import geometry;
    import graph;

    pen pen_000000 = RGB(0, 0, 0);
    pen pen_ff0000 = RGB(255, 0, 0);
    pen pen_ff7f00 = RGB(255, 127, 0);
    pen pen_0000ff = RGB(0, 0, 255);
    pen pen_0099ff = RGB(0, 153, 255);
    pen pen_9966ff = RGB(153, 102, 255);

    pair point_A = (-4.0000000000000000, 3.2200000000000000);
    pair point_B = (-6.7200000000000000, -4.2600000000000000);
    pair point_C = (3.5000000000000000, -4.3000000000000000);
    pair point_M = (-5.3599999999999990, -0.5199999999999998);
    pair point_N = (-0.2500000000000000, -0.5399999999999998);
    pair point_D = (-4.8199275125377050, 0.9651993405213131);
    pair point_E = (-2.7303162200442230, 1.9469303966310076);
    pair point_R = (-5.7699637562688530, -1.6474003297393434);
    pair point_S = (0.3848418899778885, -1.1765348016844960);
    pair point_P = (-4.4099637562688520, 2.0925996702606566);
    pair point_Q = (-3.3651581100221115, 2.5834651983155040);
    pair point_K = (-5.3894065746184920, -2.1899359966695497);
    pair point_L = (-0.0871117911968039, -2.0299250974333476);
    pair point_O = (-2.8028115577637602, 0.0291469913591767);
    pair point_O_1 = (-1.6056231155275191, -3.1617060172816460);
    path seg_f = point_A--point_B;
    path seg_g = point_B--point_C;
    path seg_h = point_A--point_C;
    path seg_i = point_B--point_E;
    path seg_j = point_D--point_C;
    path seg_l = point_M--point_O;
    path seg_m = point_O--point_N;
    path seg_q = point_B--point_O_1;
    path seg_r = point_O_1--point_C;
    path seg_s = point_O_1--point_A;
    path circ_c = circumcircle(point_B, point_D, point_C);
    path circ_d = circumcircle(point_A, point_R, point_S);
    path circ_e = circumcircle(point_P, point_E, point_N);
    path circ_k = circumcircle(point_Q, point_D, point_M);

    draw(seg_f, pen_000000);
    draw(seg_g, pen_000000);
    draw(seg_h, pen_000000);
    draw(seg_i, pen_000000);
    draw(seg_j, pen_000000);
    draw(seg_l, pen_0099ff);
    draw(seg_m, pen_0099ff);
    draw(seg_q, pen_0099ff, StickIntervalMarker(1, 2, p=pen_0099ff, size=0.3cm));
    draw(seg_r, pen_0099ff, StickIntervalMarker(1, 2, p=pen_0099ff, size=0.3cm));
    draw(seg_s, pen_9966ff+dashed);
    draw(circ_c, pen_ff0000);
    draw(circ_d, pen_ff7f00);
    draw(circ_e, pen_0000ff+dotted);
    draw(circ_k, pen_0000ff+dotted);

    dot(point_A);
    dot(point_B);
    dot(point_C);
    dot(point_M);
    dot(point_N);
    dot(point_D);
    dot(point_E);
    dot(point_R);
    dot(point_S);
    dot(point_P);
    dot(point_Q);
    dot(point_K);
    dot(point_L);
    dot(point_O);
    dot(point_O_1);

    label("$A$", point_A, dir(point_A));
    label("$B$", point_B, dir(point_B));
    label("$C$", point_C, dir(point_C));
    label("$M$", point_M, dir(point_M));
    label("$N$", point_N, 0.8NE);
    label("$D$", point_D, dir(point_D));
    label("$E$", point_E, 0.8NE);
    label("$R$", point_R, dir(point_R));
    label("$S$", point_S, E+0.5N);
    label("$P$", point_P, dir(point_P));
    label("$Q$", point_Q, 0.8NE);
    label("$K$", point_K, 1.5S);
    label("$L$", point_L, dir(point_L));
    label("$O$", point_O, NE);
    label("$O_{1}$", point_O_1, S);
\end{asy}
\end{figure}
Extends $CL$ and $BK$ intersect $AC$ and $AB$ at $E$ and $D$, respectively. It is easy to see that $BDEC$ is a cyclic quadrilateral. Notice that
\begin{align*}
&\measuredangle{LNE}=\measuredangle{LNC}=\measuredangle{LBK}=\measuredangle{LBE}\\
&\measuredangle{DCK}=\measuredangle{LCK}=\measuredangle{BMK}=\measuredangle{DMK}
\end{align*}Therefore, $ENLB$ and $DMKC$ are cyclic quadrilaterals. Denote $P$ and $Q$ by the midpoints of $AD$ and $AE$, respectively, and denote $R$ and $S$ by the midpoints of $BD$ and $CE$, respectively. For the sake of convenience, let $\Omega_A,\Omega_B$ and $\Omega_C$ be $(ARS),(ENLB)$ and $(DMKC)$, respectively. We begin with these main claims.
\begin{claim}
$P\in \Omega_B$ and $Q\in \Omega_C$.
\end{claim}
\begin{proof}
    Notice that $AE\cdot AC=AD\cdot AB$, we have
$$AE\cdot AN=\frac{1}{2}AE\cdot AC=\frac{1}{2}AD\cdot AB=AP\cdot AB.$$Therefore, $P\in \Omega_B$. Similarly, $Q\in \Omega_C$ as desired.\hfill$\square$
\end{proof}

\begin{claim}
     $K,L\in \Omega_A$.
\end{claim}
\begin{proof}
    Let $K_1$ and $L_1$ be an intersection of $\Omega_A$ with $\Omega_C$ and $\Omega_B$, respectively, so that $K_1$ and $L_1$ lies inside triangle $BMC$ and $BNC$ respectively. Notice that
\begin{align*}
    \frac{1}{2}\operatorname{Pow}(C,\Omega_A,\Omega_B)&=\operatorname{Pow}(N,\Omega_A,\Omega_B)-\frac{1}{2}\operatorname{Pow}(A,\Omega_A,\Omega_B)\\
    &=-NS\cdot NA-\frac{1}{2}(0-AE\cdot AN)\\
    &=-NS\cdot NA+\frac{1}{2}AE\cdot AN\\
    &=0
\end{align*}because $AQ=NS$, and observe further that
$$
    \operatorname{Pow}(D,\Omega_A)=-DA\cdot DR=-DP\cdot DB=\operatorname{Pow}(D,\Omega_B),
$$and $\operatorname{Pow}(L_1,\Omega_A,\Omega_B)=0$. So, $L_1,D$, and $C$ lies on the radical axis of $\Omega_A$ and $\Omega_B$; that implies $C,L_1$, and $D$ are collinear. Thus, $L_1=L$. Hence, $L\in \Omega_A$. Analogously, $K\in \Omega_A$ as desired.\hfill$\square$
\end{proof}

\begin{claim}
    The center of $(ARKLS)$ lies on the perpendicular bisector of $MN$.
\end{claim}
\begin{proof}
    Let $O_1$ be an antipode of $A$ on $(AKL)$. Notice that $RO_1$ and $SO_1$ are perpendicular to $AB$ and $AC$, respectively. Since $R$ and $S$ are midpoints of $BD$ and $CE$, we must have $O_1$ as the circumcenter of the quadrilateral $BDEC$; that is $O_1B=O_1C$. By observing the homothety of ratio $1:2$ that sends $\triangle{MON}\mapsto \triangle{BO_1C}$, we have $OM=ON$ as desired. 
    
    \hfill$\square$
\end{proof}
\noindent
From the claims above, it is enough to conclude that $OM=ON$ as desired.
\newpage
\subsubsection{Solution 2}
\begin{figure}[H]
    \centering
    \begin{asy}
        import geometry;
        size(14cm);
        defaultpen(fontsize(10pt));

        pair A = dir(100);
        pair B = dir(210);
        pair C = dir(330);
        pair M = midpoint(A--B);
        pair N = midpoint(A--C);

        pair X = A + 0.7*(B-A);
        pair Y = A + abs(X-A)*abs(B-A)/abs(C-A)^2*(C-A);   // AX*AB = AY*AC forces (BXYC) cyclic
        pair R = circumcenter(B, X, C);
        pair O = midpoint(A--R);

        path lineBY = (B-0.2*(Y-B))--(B+2.0*(Y-B));
        path lineCX = (C-0.2*(X-C))--(C+2.0*(X-C));
        pair[] kk = intersectionpoints(circumcircle(C, X, M), lineBY);
        pair K  = (kk[0].y < kk[1].y) ? kk[0] : kk[1];
        pair Kp = (kk[0].y < kk[1].y) ? kk[1] : kk[0];
        pair[] ll = intersectionpoints(circumcircle(B, Y, N), lineCX);
        pair L  = (ll[0].y < ll[1].y) ? ll[0] : ll[1];
        pair Lp = (ll[0].y < ll[1].y) ? ll[1] : ll[0];
        pair D = extension(B, Y, C, X);
        pair Z = midpoint(B--X);
        pair W = midpoint(C--Y);
        pair Xp = midpoint(A--X);

        draw(circumcircle(B, X, C), red);
        draw(circumcircle(C, X, M), blue);
        draw(circumcircle(B, Y, N), blue);
        draw(circle(O, abs(O-A)), dashed+deepgreen+0.8);

        draw(A--B--C--cycle, black+1);
        draw(B--Kp, black+0.6);
        draw(C--Lp, black+0.6);
        draw(M--N, dotted+gray(0.5));
        draw(A--R, dotted+gray(0.5));
        draw(O--M, dashed+purple);
        draw(O--N, dashed+purple);

        dot("$A$", A, dir(100));
        dot("$B$", B, dir(190));
        dot("$C$", C, dir(-25));
        dot("$M$", M, dir(145));
        dot("$N$", N, 1.6*dir(42));
        dot("$X$", X, dir(110));
        dot("$Y$", Y, 1.3*dir(85));
        dot("$K$", K, 1.5*dir(-115));
        dot("$L$", L, dir(-75));
        dot("$K'$", Kp, dir(50));
        dot("$L'$", Lp, dir(170));
        dot("$D$", D, dir(-87));
        dot("$Z$", Z, 1.5*dir(145));
        dot("$W$", W, dir(43));
        dot("$X'$", Xp, dir(145));
        dot("$R$", R, dir(-85));
        dot("$O$", O, 1.2*dir(60));
    \end{asy}
\end{figure}
\noindent
Let $D = BK \cap CL$, $X = CL \cap AB, Y = BL \cap AC$ and let $Z, W$ denote midpoints of $BX, CY$, respectively.

\medskip
It's easy to note that $(BXYC)$ are cyclic by the given angle condition. The key claim is to see that $(AZWKL)$ are cyclic, which proves the problem as for the center $R$ of $\odot(BXYC)$, the midpoint of $AR$ is the circumcenter of $AZW$ which trivially lies on the perpendicular bisector of $MN$.

\medskip
Hence, it suffices to prove that $(AZWKL)$ are cyclic.

\medskip
To see this, first note that by the given angle conditions, it's easy to prove $(BYNL)$ and $(CXMK)$ are cyclic respectively by angle chasing. Construct $K'$ as the second intersection of $\odot(CXMK)$ and line $BY$, and define $L'$ similarly. Then we observe the following:
\begin{itemize}
    \item By PoP on $D$,$$DL \cdot DL' = DB \cdot DY = DX \cdot DC = DK \cdot DK'.$$Hence $(KK'LL')$ are cyclic.
    \item Let $X'$ be the midpoint of $AX$. Then$$AB \cdot AX' = \frac{1}{2} AB \cdot AC = AY \cdot AN,$$so $X'$ lies on $\odot (BLNYL')$.
    \item By PoP at $X$,$$AX \cdot XZ = AX' \cdot XB = XL \cdot XL'.$$Hence $(AZLL')$ are cyclic. Similarly, $(AWKK')$ are cyclic.
    \item Finally, by PoP at $B$, we have$$BK \cdot BK' = BM \cdot BX = BA \cdot BZ,$$meaning that $(AKK'Z)$ are cyclic, as well as $(ALL'W)$ by symmetry.
\end{itemize}
\noindent
This proves the result. We're done.
\newpage
\subsubsection{Solution 3}
\begin{figure}[H]
    \centering
    \begin{tikzpicture}[scale=7] % scale=7 對應到直徑約 14cm 的圖形大小

    % 1. 定義初始點 A, B, C (對應 Asymptote 中的 dir(angle))
    \tkzDefPoint(130:1){A}
    \tkzDefPoint(210:1){B}
    \tkzDefPoint(-30:1){C}

    % 2. 尋找 E 點: E = 0.3*C + 0.7*A
    % 透過重心座標 (barycentric coordinates) 定義比例
    \tkzDefBarycentricPoint(C=0.3,A=0.7) \tkzGetPoint{E}

    % 3. B, C, E 的外接圓與圓心 O
    \tkzDefCircle[circum](B,C,E) \tkzGetPoint{O}

    % 4. 尋找 F 點: circumcircle(B,C,E) 與 A--B 的交點
    % 幾何技巧: 因為 B 已經在圓上且在線段 AB 上，
    % F 點必定是 B 點對「O 點至 AB 投影點」的對稱點，這樣可避免交點選擇的模糊性。
    \tkzDefPointBy[projection=onto A--B](O) \tkzGetPoint{projO}
    \tkzDefPointBy[reflection=over O--projO](B) \tkzGetPoint{F}

    % 5. 定義中點 M, N, M', N'
    \tkzDefMidPoint(A,B) \tkzGetPoint{M}
    \tkzDefMidPoint(A,C) \tkzGetPoint{N}
    \tkzDefMidPoint(F,B) \tkzGetPoint{M1}
    \tkzDefMidPoint(E,C) \tkzGetPoint{N1}

    % 6. 尋找外接圓圓心 O_MFC 與 O_NEB
    \tkzDefCircle[circum](M,F,C) \tkzGetPoint{OMFC}
    \tkzDefCircle[circum](N,E,B) \tkzGetPoint{ONEB}

    % 7. 尋找 K 點: circumcircle(M,F,C) 與 B--E 的交點
    \tkzInterLC(B,E)(OMFC,C) \tkzGetPoints{K1}{K2}
    \coordinate (K) at (K1); % 若圖形上 K 跑到線段外，可將此處 K1 改為 K2

    % 8. 尋找 L 點: circumcircle(N,E,B) 與 C--F 的交點
    \tkzInterLC(C,F)(ONEB,B) \tkzGetPoints{L1}{L2}
    \coordinate (L) at (L1); % 若圖形上 L 跑到線段外，可將此處 L1 改為 L2

    % 9. 尋找 A, K, L 的外接圓圓心
    \tkzDefCircle[circum](A,K,L) \tkzGetPoint{OAKL}

    % ================= 繪製階段 =================
    
    % 繪製線段
    \tkzDrawSegment[color=red](B,E)
    \tkzDrawSegment[color=red](C,F)
    \tkzDrawPolygon(A,B,C)

    % 繪製外接圓
    \tkzDrawCircle[color=red](O,B)
    \tkzDrawCircle[color=blue](OMFC,C)
    \tkzDrawCircle[color=blue](ONEB,B)
    \tkzDrawCircle[color=blue](OAKL,A)

    % 畫出所有點
    \tkzDrawPoints(A,B,C,E,F,M,N,M1,N1,O,K,L)

    % 標記點位文字 (依照原 Asymptote dir 函數推斷的最佳文字方位)
    \tkzLabelPoint[above left](A){$A$}
    \tkzLabelPoint[below left](B){$B$}
    \tkzLabelPoint[below right](C){$C$}
    \tkzLabelPoint[above right](E){$E$}
    \tkzLabelPoint[above left](F){$F$}
    \tkzLabelPoint[left](M){$M$}
    \tkzLabelPoint[right](N){$N$}
    \tkzLabelPoint[left](M1){$M'$}
    \tkzLabelPoint[right](N1){$N'$}
    \tkzLabelPoint[below](O){$O$}
    \tkzLabelPoint[above right](K){$K$}
    \tkzLabelPoint[above left](L){$L$}

\end{tikzpicture}
\end{figure}
Let $E$ and $F$ be $BK \cap AC$ and $CL \cap AB$. Note that $\triangle AEB \sim \triangle AFC$ so $BCEF$ is cyclic. Let $O$ be its circumcenter, and $N'$ and $M'$ the midpoints of $CE$ and $BF$ respectively.

\medskip
The angle condition is equivalent to $CKMF$ and $BLNE$ being cyclic. By Reim's $MEFN$ is cyclic too. I claim that $(AO)$ and $(CMF)$ have radical axis $BE$. Indeed,
\[BM' \cdot BA = \frac 12 BF\cdot BA = BM \cdot BF\]
Meanwhile, by Linearity of Power of a Point, if $$f(\bullet) = \mathrm{Pow}_{(MFC)}(\bullet)-\mathrm{Pow}_{(AN'M')}(\bullet),$$ considering $f(E)$, we need to show
\[0= 0 \cdot AC = AE f(C) + ECf(A) = - AE\cdot CN' \cdot CA + EC \cdot AF \cdot AM.  \]This is equivalent to
\[AE\cdot AN \cdot EC = AE\cdot CA\cdot CN'.\]Yet this is true as $$AN \cdot EC = \frac 12 AC \cdot EC  = CA \cdot CN'.$$

\medskip
Analogously, $C$ lies on the radical axis of $(AO)$ and $(BNE)$. Hence $K$ and $L$ are simply two of the intersection points between $(AN'M')$ and $(MFC)$ and $(NEB)$ respectively. So $(AKL)=(AO)$ has $A$-antipode $O$ and dilating by half, it only remains to show that the midpoint of $AO$ lies on the perpendicular bisector of $MN$. This is evident as the image of $(BCEF)$ has a base $MN$ and the definition of circumcenter suffices.

\newpage
\subsubsection{Solution 4}
\begin{figure}[H]
    \centering
    \begin{asy}
        import graph; 
        size(14cm); 
        real labelscalefactor = 0.5; /* changes label-to-point distance */
        pen dps = linewidth(0.2) + fontsize(10); defaultpen(dps); /* default pen style */ 
        pen dotstyle = black; /* point style */ 
        real xmin = -7.8, xmax = 11.48, ymin = -7.37, ymax = 5.57;  /* image dimensions */
        pen wewdxt = rgb(0.43137254901960786,0.42745098039215684,0.45098039215686275); 
        pen bqqqsq = rgb(0.6901960784313725,0,0.12549019607843137); 
        pen rutezs = rgb(0.0784313725490196,0.24313725490196078,0.5725490196078431); 

        /* define paths for intersections later */
        path lineAB_path = (-3.88049,4.89775)--(-5.84,-2.92);
        path lineAC_path = (-3.88049,4.89775)--(3.68,-3.4);
        path w2_path = circle((0.23856,0.40695), 5.13192);
        path w1_path = circle((-3.66286,-0.00182), 3.64086);

        /* draw figures */
        draw(lineAB_path, linewidth(1) + wewdxt); 
        draw((-5.84,-2.92)--(3.68,-3.4), linewidth(1) + wewdxt); 
        draw(lineAC_path, linewidth(1) + wewdxt); 
        draw(circle((-0.98937,-1.36250), 5.09455), linewidth(0.8) + linetype("2 2") + red); 
        draw((-2.58486,3.47577)--(-5.84,-2.92), linewidth(1) + wewdxt); 
        draw((-4.53193,2.29875)--(3.68,-3.4), linewidth(1) + wewdxt); 
        draw(w2_path, linewidth(0.8) + linetype("2 2") + bqqqsq); 
        draw(w1_path, linewidth(0.8) + linetype("2 2") + bqqqsq); 
        draw(circle((-2.43493,1.76763), 3.44773), linewidth(1) + rutezs); 
        draw((-5.18596,-0.31063)--(-0.98937,-1.36250), linewidth(0.8) + linetype("2 2") + wewdxt); 
        draw((-0.98937,-1.36250)--(0.54757,0.03789), linewidth(0.8) + linetype("2 2") + wewdxt); 

        /* dots and labels */
        dot((-3.88049,4.89775),dotstyle); 
        label("$A$", (-4.10,5.00), NE * labelscalefactor); 
        dot((-5.84,-2.92),dotstyle); 
        label("$B$", (-6.26,-3.23), NE * labelscalefactor); 
        dot((3.68,-3.4),dotstyle); 
        label("$C$", (3.76,-3.51), NE * labelscalefactor); 
        dot((-4.86025,0.98887),linewidth(4pt) + dotstyle); 
        label("$M$", (-4.77,1.16), NE * labelscalefactor); 
        dot((-0.10025,0.74887),linewidth(4pt) + dotstyle); 
        label("$N$", (-0.01,0.92), NE * labelscalefactor); 
        dot((-4.53193,2.29875),linewidth(4pt) + dotstyle); 
        label("$L'$", (-5.0,2.20), NE * labelscalefactor); 
        dot((-2.58486,3.47577),linewidth(4pt) + dotstyle); 
        label("$K'$", (-2.61,3.65), NE * labelscalefactor); 
        dot((-4.75313,-0.78449),linewidth(4pt) + dotstyle); 
        label("$K$", (-4.67,-0.9), NE * labelscalefactor); 
        dot((-0.10474,-0.77354),linewidth(4pt) + dotstyle); 
        label("$L$", (-0.02,-0.92), NE * labelscalefactor); 
        dot((-251.29229,114.24387),linewidth(4pt) + dotstyle); 
        label("$W$", (-251.21,114.42), NE * labelscalefactor); 
        dot((-0.98937,-1.36250),linewidth(4pt) + dotstyle); 
        label("$A'$", (-0.91,-1.60), NE * labelscalefactor); 
        dot((0.54757,0.03789),linewidth(4pt) + dotstyle); 
        label("$F$", (0.62,0.0), NE * labelscalefactor); 
        dot((-5.18596,-0.31063),linewidth(4pt) + dotstyle); 
        label("$G$", (-5.70,-0.45), NE * labelscalefactor); 

        /* circle labels */
        label("$\Omega$", (1.18937 + 5.09455*dir(-75)), SE, red); 
        label("$\omega_2$", (0.45856 + 5.13192*dir(20)), NE, bqqqsq); 
        label("$\omega_1$", (-3.66286 + 3.64086*dir(160)), NW, bqqqsq); 
        label("$\Gamma$", (-2.83493 + 3.44773*dir(135)), NW, rutezs); 

        /* New intersection points */
        pair[] int_AK = intersectionpoints(lineAC_path, w2_path);
        pair[] int_AL = intersectionpoints(lineAB_path, w1_path);

        /* Plot intersection of AK' and w2 (ignoring C) */
        for(int i=0; i<int_AK.length; ++i) {
            if(length(int_AK[i] - (3.68,-3.4)) > 0.5) {
                dot(int_AK[i], linewidth(4pt) + dotstyle);
                label("$X$", int_AK[i], N * labelscalefactor * 3);
            }
        }

        /* Plot intersection of AL' and w1 (ignoring B) */
        for(int i=0; i<int_AL.length; ++i) {
            if(length(int_AL[i] - (-5.84,-2.92)) > 0.5) {
                dot(int_AL[i], linewidth(4pt) + dotstyle);
                label("$Y$", int_AL[i], NW * labelscalefactor);
            }
        }

        clip((xmin,ymin)--(xmin,ymax)--(xmax,ymax)--(xmax,ymin)--cycle); 
    \end{asy}
\end{figure}
Let $K'= BK \cap AC$ and $L' = CL \cap AB$. From the conditions of the problem, we have that $BLNK'$, $CKML'$ and $BL'K'C$ are respectively cyclic. Call these circles $\omega_1, \omega_2$ and $\Omega$ for convenience.

Let $X, Y$ be intersection of $\omega_2$ and $\omega_1$ with $AC$ and $AB$ respectively. We claim that $X$ is the midpoint of $AK’$ and analogously $Y$ is the midpoint of $AL'$. Indeed, just note that
\[ AX \cdot AC = \mathrm{Pow}_A(\omega_2) = AL' \cdot AM = \frac{1}{2} AL' \cdot AB \implies AX =  \frac{1}{2} AL' \cdot \frac{AB}{AC} = \frac{1}{2} AK' \]since $\triangle AK'L' \sim \triangle ABC$.

Let $A'$ be the $A$-antipode of $(AKL)$. Therefore, to prove that $OM = ON$, it suffices to prove that $A'B = A'C$ by homothety $\mathcal{H}(A, 2)$. We will use phantom point to do this: we instead assume that $A'$ is the circumcenter of $\Omega$, and we will show that $K$ and $L$ lies on circle with diameter $AA'$, which we will call $\Gamma$ for convenience.

To do this, we will show that $BK' \equiv \mathrm{rad}(\Gamma, \omega_2)$. Indeed, let projection of $A'$ to $AB$ and $AC$ be $G$ and $F$ respectively. Then $G, F \in \Gamma$. We also have the following two power calculation, which proves $BK' \equiv \mathrm{rad}(\Gamma, \omega_2)$:
\begin{align*} 
\mathrm{Pow}_B(\Gamma) &= BG \cdot BA = \frac{1}{2} BL' \cdot BA = BL' \cdot BM = \mathrm{Pow}_{B}(\omega_2) \\
\mathrm{Pow}_{K'}(\Gamma) &= K'A \cdot K'F = \frac{1}{2} K'A \cdot K'C = K'X \cdot K'C = \mathrm{Pow}_{K'}(\omega_2)
\end{align*}This proves $BK' \equiv \mathrm{rad}(\Gamma, \omega_2)$, which means that $K \in BK' \cap \omega_2$ must lie on $\Gamma$. Analogously, we can show that $L$ lies on $\Gamma$ as well, which proves what we wanted.
\newpage 

\subsubsection{Solution 5}
\begin{figure}[H]
    \centering
    \begin{asy}
        import geometry; 
        size(14cm); 
        real labelscalefactor = 0.5; /* changes label-to-point distance */
        pen dps = linewidth(0.7) + fontsize(10); defaultpen(dps); /* default pen style */ 
        pen dotstyle = black; /* point style */ 
        real xmin = -2.5952983660369107, xmax = 9.304851108401753, ymin = -2.890416452859565, ymax = 6.339981177716112;  /* image dimensions */
        pen fuqqzz = rgb(0.9568627450980393,0,0.6); 

        draw(arc((4.5,3),0.3591594408784305,-86.36641220702231,-50.19442890773481)--(4.5,3)--cycle, linewidth(0.8) + fuqqzz); 
        draw(arc((0,0),0.3591594408784305,20.466884611212286,56.63886791049988)--(0,0)--cycle, linewidth(0.8) + fuqqzz); 
        draw(arc((1,3),0.3591594408784305,-108.43494882292202,-88.11393179920499)--(1,3)--cycle, linewidth(0.8) + blue); 
        draw(arc((7,0),0.3591594408784305,144.73175435884332,165.05277138256034)--(7,0)--cycle, linewidth(0.8) + blue); 
        draw(arc((0,0),0.3591594408784305,71.56505117707799,91.88606820079501)--(0,0)--cycle, linewidth(0.8) + blue); 
        draw(arc((7,0),0.3591594408784305,93.63358779297764,129.8055710922652)--(7,0)--cycle, linewidth(0.8) + fuqqzz); 
         /* draw figures */
        draw((2,6)--(0,0), linewidth(0.8)); 
        draw((0,0)--(7,0), linewidth(0.8)); 
        draw((7,0)--(2,6), linewidth(0.8)); 
        draw(circle((3.5,1.0589320846622627), 3.6566839021068205), linewidth(0.8)); 
        draw((7,0)--(1.3353592507973573,4.0060777523920725), linewidth(0.8)); 
        draw((0,0)--(3.0895749986928567,4.692510001568572), linewidth(0.8)); 
        draw(circle((2.75,3.5294660423311313), 2.5818671607956523), linewidth(0.8)); 
        draw((1,3)--(1.04645194936694,1.5893727262302526), linewidth(0.8) + blue); 
        draw((7,0)--(-0.062083242746548625,1.885309801975236), linewidth(0.8) + blue); 
        draw((0,0)--(6.837933730153274,2.5520950272215157), linewidth(0.8) + fuqqzz); 
        draw((4.5,3)--(4.581913479911532,1.7100894902910002), linewidth(0.8) + fuqqzz); 
        draw((0,0)--(-0.062083242746548625,1.885309801975236), linewidth(0.8) + blue); 
        draw((7,0)--(6.837933730153274,2.5520950272215157), linewidth(0.8) + fuqqzz); 
         /* dots and labels */
        dot((2,6),linewidth(2pt) + dotstyle); 
        label("$A$", (1.9181386076686993,6.052653625013368), NE * labelscalefactor); 
        dot((0,0),linewidth(2pt) + dotstyle); 
        label("$B$", (-0.34456586986541277,-0.32841244126009045), NE * labelscalefactor); 
        dot((7,0),linewidth(2pt) + dotstyle); 
        label("$C$", (7.14989446313117,-0.3403844226227048), NE * labelscalefactor); 
        dot((1,3),linewidth(2pt) + dotstyle); 
        label("$M$", (0.70896849004465,2.939938470733632), NE * labelscalefactor); 
        dot((4.5,3),linewidth(2pt) + dotstyle); 
        label("$N$", (4.587890451531699,3.0955742284476186), NE * labelscalefactor); 
        dot((1.3353592507973573,4.0060777523920725),linewidth(2pt) + dotstyle); 
        label("$D$", (1.0681279309230804,4.125164625632454), NE * labelscalefactor); 
        dot((3.0895749986928567,4.692510001568572),linewidth(2pt) + dotstyle); 
        label("$E$", (3.1392807066553625,4.735735675125786), NE * labelscalefactor); 
        dot((3.5,1.0589320846622627),linewidth(2pt) + dotstyle); 
        label("$Z$", (3.7618237375113086,0.8807576763639606), NE * labelscalefactor); 
        dot((2.75,3.5294660423311313),linewidth(2pt) + dotstyle); 
        label("$O$", (2.792093247139547,3.574453482952193), NE * labelscalefactor); 
        dot((1.04645194936694,1.5893727262302526),linewidth(2pt) + dotstyle); 
        label("$K$", (1.1758757631866097,1.6828804276591232), NE * labelscalefactor); 
        dot((4.581913479911532,1.7100894902910002),linewidth(2pt) + dotstyle); 
        label("$L$", (4.192815066565426,1.6230205208460513), NE * labelscalefactor); 
        dot((-0.062083242746548625,1.885309801975236),linewidth(2pt) + dotstyle); 
        label("$F$", (-0.32062190714018407,1.9462640176366393), NE * labelscalefactor); 
        dot((6.837933730153274,2.5520950272215157),linewidth(2pt) + dotstyle); 
        label("$G$", (6.994258705417184,2.5688070484925865), NE * labelscalefactor); 
        clip((xmin,ymin)--(xmin,ymax)--(xmax,ymax)--(xmax,ymin)--cycle); 
    \end{asy}
\end{figure}
Define $D = CL \cap BA$, $E = BK \cap CA$, $F = CK \cap (DECB)$, $G = BL \cap (DECB)$, and $Z$ as the circumcenter of $(DECB)$. 

\medskip
\noindent
First, observe that $MK \parallel BF$ and $CG \parallel NL$. 
\begin{claim}
    $\angle{AKZ} = 90^{\circ}.$
\end{claim}
\begin{proof}
    We will use complex numbers and take $(BCED)$ as the unit circle. Consider that $MK \parallel BF$ implies
\[\frac{k-m}{b-f} \in \mathbb{R}\Longrightarrow \frac{2k-a-b}{b-f} = \overline{\left(\frac{2k-a-b}{b-f}\right)}  \Longrightarrow 2k-a-b-f+2bf\overline{k} - \overline{a}bf = 0.  \]
Then, since $A \in CE$, we know $a + ce\overline{a} = c + e$. Substituting to the previous equation,
\[2k - c - e - b - f + ce\overline{a} + 2bf\overline{k} - \overline{a}bf = 0. \]
Furthermore, $K \in CF, BE$, hence $k + cf\overline{k} = c + f$ and $k+be\overline{k} = b + e$. Substituting,
\[2k - 2k - cf\overline{k} - be\overline{k} + ce\overline{a} + 2bf\overline{k} - \overline{a}bf = 0 \Longrightarrow  cf\overline{k} - be\overline{k} + ce\overline{a} + 2bf\overline{k} - \overline{a}bf = 0. \]Then, we get
\[\frac{\overline{a}}{\overline{k}} = \frac{2bf-cf-be}{bf-ce}. \]
Then, we simply need to prove that $\frac{a-k}{k-z} = \frac{a}{k} - 1$ is imaginary which is simply some trivial computations. \hfill$\square$
\end{proof}
\noindent
Analogously, $\angle{ALZ} = 90^{\circ}$. Therefore, $Z$ is the antipode of $A$ on $(AKL)$. Hence, by homothety on center $A$ with scale $\frac{1}{2}$, then $B \rightarrow M$, $C \rightarrow N$, $Z \rightarrow O$. Since $ZB = ZC$, then $OM = ON$. \hfill$\square$
\newpage
\subsubsection{Solution 6}
Barycentric coordinates on $\triangle ABC$. Let $K=(x_1, y_1, z_1)$ and $L=(x_2, y_2, z_2)$. Let line $BK$ meet $AC$ at $D=(x_1:0:z_1)$, and let line $CL$ meet $AB$ at $E=(x_1:y_1:0)$. Note that the first angle condition implies $\triangle ABD \sim \triangle ACE$, so
\[
\frac{AD}{AE} = \frac{AB}{AC} \implies \frac{c \cdot \frac{z_1}{x_1+z_1}}{b \cdot \frac{y_2}{x_2+y_2}} = \frac{c}{b}.
\]
Thus,
\[
\boxed{b^2z_1(x_2+y_2) = c^2y_2(x_1+z_1)}.
\]

\noindent Let $X$ be the point on line $BK$ such that $\overleftrightarrow{XL} \parallel \overleftrightarrow{AB}$. We can easily compute $X$:
\[
X = (x_2,y_2,z_2) + \left( \frac{x_1z_2-x_2z_1}{z_1}, -\frac{x_1z_2-x_2z_1}{z_1}, 0 \right) = \left( \frac{x_1z_2}{z_1}, y_2 - \frac{x_1z_2-x_2z_1}{z_1}, z_2 \right).
\]

\noindent Since $K$ lies inside $\angle LBA$, $\overrightarrow{LX}$ is in the same direction as $\overrightarrow{BA}$. It follows that $\frac{x_1z_2-x_2z_1}{z_1}$ is positive, so
\[
XL = c \cdot \frac{x_1z_2-x_2z_1}{z_1}.
\]
By looking at the $C$-coordinates of $\overrightarrow{BK}$ and $\overrightarrow{BX}$, we have
\[
BX = BK \cdot \frac{z_2}{z_1}.
\]

\noindent We can also compute the following areas and segments:
\begin{align*}
CN &= \frac{1}{2}b \\
\frac{[BLK]}{[ABC]} &=
\begin{vmatrix}
0 & 1 & 0 \\
x_2 & y_2 & z_2 \\
x_1 & y_1 & z_1
\end{vmatrix}
= x_1z_2 - x_2z_1 \\
\frac{[CNL]}{[ABC]} &=
\begin{vmatrix}
0 & 0 & 1 \\
\frac{1}{2} & 0 & \frac{1}{2} \\
x_2 & y_2 & z_2
\end{vmatrix}
= \frac{1}{2}y_2.
\end{align*}

\noindent Note that $\angle LXB = \angle ABK = \angle LCN$ and $\angle LBX = \angle LNC$, so $\triangle LNC \sim \triangle LBX$. Hence,
\[
\frac{XL}{CL} = \frac{BL}{NL} = \frac{BX}{CN} = BK \cdot \frac{z_2}{z_1} \cdot \frac{2}{b}. \tag{1}
\]
Additionally, we have
\[
\frac{[BLK]}{BK \cdot BL} = \frac{1}{2}\sin\angle LBK = \frac{1}{2}\sin\angle LNC = \frac{[CNL]}{CN \cdot NL}. \tag{2}
\]

\noindent From (1), we get
\[
BK \cdot CL = \frac{b}{2} \cdot \frac{z_1}{z_2} \cdot XL = \frac{bc}{2} \cdot \frac{x_1z_2-x_2z_1}{z_2}.
\]
By symmetry, it follows that
\[
\boxed{\frac{x_1z_2-x_2z_1}{z_2} = \frac{x_2y_1-x_1y_2}{y_1}}.
\]
From (2), we get
\begin{align*}
\frac{BK}{CN} \cdot \frac{BL}{NL} &= \frac{[BLK]}{[CNL]} \\
\implies \frac{BK}{\frac{1}{2}b} \cdot BK \cdot \frac{z_2}{z_1} \cdot \frac{2}{b} &= \frac{x_1z_2-x_2z_1}{\frac{1}{2}y_2},
\end{align*}
so
\[
\boxed{BK^2 = \frac{b^2}{2} \cdot \frac{z_1}{y_2} \cdot \frac{x_1z_2-x_2z_1}{z_2}}.
\]

\noindent Suppose $(AKL)$ is given by
\[
-a^2yz - b^2zx - c^2xy + (x+y+z)(vy+wz) = 0.
\]
Then,
\begin{align*}
y_1v + z_1w &= -P_K \\
y_2v + z_2w &= -P_L,
\end{align*}
where $P_K := \mathrm{Pow}_{(ABC)}(K)$ and $P_L := \mathrm{Pow}_{(ABC)}(L)$. By Cramer's rule, we have
\begin{align*}
v &= \frac{
\begin{vmatrix}
-P_K & z_1 \\
-P_L & z_2
\end{vmatrix}
}{y_1z_2-y_2z_1}, \\
w &= \frac{
\begin{vmatrix}
y_1 & -P_K \\
y_2 & -P_L
\end{vmatrix}
}{y_1z_2-y_2z_1}.
\end{align*}

\noindent It suffices to prove that $\mathrm{Pow}_{(AKL)}(M) = \mathrm{Pow}_{(AKL)}(N)$, or
\[
-\frac{1}{4}c^2 + \frac{1}{2}v = -\frac{1}{4}b^2 + \frac{1}{2}w \implies 2
\begin{vmatrix}
P_K & y_1+z_1 \\
P_L & y_2+z_2
\end{vmatrix}
= (b^2-c^2)(y_1z_2-y_2z_1). \tag{3}
\]

\noindent Note that
\begin{align*}
BK^2 &= P_K + a^2z_1 + c^2x_1 \\
\implies P_K &= \frac{b^2}{2} \cdot \frac{z_1}{y_2} \cdot \frac{x_1z_2-x_2z_1}{z_2} - a^2z_1 - c^2x_1.
\end{align*}

\noindent We will get rid of $a$. Note that
\[
2y_2z_2BK^2 = 2y_2z_2 \Big(a^2z_1(x_1+z_1) - b^2x_1z_1 + c^2x_1(x_1+z_1)\Big) = b^2z_1(x_1+y_1+z_1)(x_1z_2-x_2z_1),
\]
so
\[
2a^2y_2z_1z_2(x_1+z_1) = 2b^2x_1y_2z_1z_2 + b^2z_1(x_1+y_1+z_1)(x_1z_2-x_2z_1) - 2c^2x_1y_2z_2(x_1+z_1).
\]

\noindent Now,
\begin{align*}
2y_2z_2P_K &= b^2z_1(x_1z_2-x_2z_1) - \frac{2b^2x_1y_2z_1z_2 + b^2z_1(x_1+y_1+z_1)(x_1z_2-x_2z_1)}{x_1+z_1} \\
&\quad + 2c^2x_1y_2z_2 - 2c^2x_1y_2z_2 \\
&= -b^2z_1 \cdot \frac{2b^2x_1y_2z_2 + y_1(x_1z_2-x_2z_1)}{x_1+z_1} \\
&= -b^2z_1 \cdot \frac{2x_1y_2z_2 + z_2(x_2y_1-x_1y_2)}{x_1+z_1} \\
&= -b^2z_1z_2 \cdot \frac{x_1y_2+x_2y_1}{x_1+z_1}.
\end{align*}

\noindent There is a similar equation for $P_L$. Multiplying (3) by $y_2z_1$ gives
\[
-b^2z_1^2 \cdot \frac{x_1y_2+x_2y_1}{x_1+z_1}(y_2+z_2) + c^2y_2^2 \cdot \frac{x_2z_1+x_1z_2}{x_2+y_2}(y_1+z_1) = (b^2-c^2)y_2z_1(y_1z_2-y_2z_1).
\]

\noindent Multiplying this by $(x_1+z_1)(x_2+y_2)$, moving the terms with $b^2$ to the RHS, and moving the terms with $c^2$ to the LHS gives
\begin{align*}
\mathrm{RHS} &= b^2z_1(x_2+y_2) \Big( y_2(y_1z_2-y_2z_1)(x_1+z_1) + z_1(x_1y_2+x_2y_1)(y_2+z_2) \Big) \\
&= b^2z_1(x_2+y_2) \Big( x_1y_1y_2z_2 + y_1y_2z_1z_2 - x_1y_2^2z_1 - y_2^2z_1^2 + x_1y_2^2z_1 + x_1y_2z_1z_2 + x_2y_1y_2z_1 + x_2y_1z_1z_2 \Big) \\
&= b^2z_1(x_2+y_2) \Big( y_1y_2z_1z_2 - y_2^2z_1^2 + (x_1y_1y_2z_2 + x_2y_1z_1z_2) + (x_1y_2z_1z_2 + x_2y_1y_2z_1) \Big).
\end{align*}

\noindent Since the term
\[
y_1y_2z_1z_2 - y_2^2z_1^2 + (x_1y_1y_2z_2 + x_2y_1z_1z_2) + (x_1y_2z_1z_2 + x_2y_1y_2z_1)
\]
is symmetric in $B$ and $C$, we similarly have
\[
\mathrm{LHS} = c^2y_2(x_1+z_1) \Big( y_1y_2z_1z_2 - y_2^2z_1^2 + (x_1y_1y_2z_2 + x_2y_1z_1z_2) + (x_1y_2z_1z_2 + x_2y_1y_2z_1) \Big).
\]

\noindent Since $b^2z_1(x_2+y_2) = c^2y_2(x_1+z_1)$, it follows that $\mathrm{LHS} = \mathrm{RHS}$. 
\hfill$\square$
\newpage
\subsection{International Mathematical Olympiad/Problem 3}
\begin{problem*}
    Let $n$ be a positive integer. Liu Bang and Xiang Yu have a stick of length 1 and want to divide it between themselves. Liu marks at most $n$ points on the stick, and then Xiang marks at most $n$ points on the stick. The marked points are distinct. Then, the stick is cut at all marked points, creating a number of pieces. Afterwards, they take turns claiming any unclaimed piece of the stick, with Liu going first. Each player's goal is to maximise the total length of their own pieces.

\medskip
\noindent
For each $n$, determine the largest value $c$ such that Liu may guarantee a total length of at least $c$, regardless of Xiang's play.
\end{problem*}
\subsubsection{Solution 1}
Suppose that Liu Bang cuts the stick into pieces of lengths $a_1,\dots,a_{n+1}$. Xiang Yu needs to make $n$ more cuts so that the Liu Bang gets at most $\tfrac{2^n}{2^{n+1}-1}$.
\begin{claim}
    There exists disjoint subsets $I$ and $J$ such that
$$0\leq \sum_{i\in I} a_i - \sum_{j\in J} a_j \leq \frac 1{2^{n+1}-1}.$$
\end{claim}
\begin{proof}
    Pigeonhole argument. There are $2^{n+1}$ subset sums of the form $\sum_{i\in I}a_i$, so by pigeonhole two of those differ by at most $\tfrac 1{2^{n+1}-1}$. By cancelling terms appearing in both, we can make $I$ and $J$ disjoint. \hfill$\square$
\end{proof}
\noindent
Let $$d=\sum_{i\in I}a_i - \sum_{j\in J}a_j.$$ We take $(I,J)$ so that $d$ is minimal. Note that if there is $i\in I$ such that $a_i\leq d$, then we can remove $i$ from $I$ and get a subset with smaller difference. Thus, we assume that $a_i>d$ for all $i$.

\medskip
Now, we place the following sticks along an arc of a circle in that order:
\begin{itemize}
    \item a dummy stick of length $d$;
    \item the sticks $a_i$ for $i\in I$; and
    \item the sticks $a_j$ for $j\in J$.
\end{itemize}
Now, Xiang Yu creates the cuts as follow: reflect every point between each stick across the center of the circle. Since two of them are already diametrically opposite, this introduces at most $|I|+|J|-1$ points. Cut the sticks not in $I$ or $J$ into two pieces of equal lengths. The upshot of this is, by symmetry, each length appears even number of times, except $d$. Thus, the difference between what Liu Bang and Xiang Yu get is at most $d$.

\medskip
Here is the illustration where the sticks in $I$ have lengths $8,9$, while sticks in $J$ have lengths $2,4,5$. Thus, $d=6$. The dummy stick is in black, the sticks in $I$ are in red, while the sticks in $J$ are in blue.

\medskip
\begin{figure}[H]
    \centering
    \begin{asy}
        size(14cm,0);
        int n = 34;
        int[] p = {6,8,9,2,4,5};
        int[] q = {0};
        int s=0;
        for(int i=0; i<p.length; ++i){
        s += p[i];
        q.push(s);
        }
        for(int i=0; i<p.length; ++i){
        pen s = black+linewidth(1);
        if(i>=1 && i<=2){
        s = red + linewidth(1);
        }
        else if(i>=3){
        s = blue + linewidth(1);
        }
        draw(arc((0,0),1,q[i]*360/n+3,q[i+1]*360/n-3,CCW), s);
        real x = (q[i]+q[i+1])/2*360/n;
        label("$" + ((string)p[i]) + "$", dir(x), dir(x), s);
        }
        for(int i=0; i<p.length; ++i){
        //if(i==1) continue;
        draw(dir(q[i]*360/n) -- (-dir(q[i]*360/n)), dashed);
        dot(-dir(q[i]*360/n));
        }
    \end{asy}
\end{figure}

\newpage
\subsubsection{Solution 2}
The answer is $$\boxed{\frac{2^n}{2^{n+1}-1}}.$$ Say Liu Bang's moves cut (not mark) the stick into blocks, which are then cut (not marked) into slices by Xiang Yu. Note that the second part of the game is deterministic: each player simply picks the largest remaining slice. We allow cuts to overlap, producing blocks and slices of length $0$. For now, we will WLOG force Liu Bang to make exactly $n$ cuts, but not Xiang Yu; this is solely to make the presentation of the solution more intuitive.
\paragraph{\textcolor{blue}{Strategy for Xiang Yu:}} Let the lengths of Liu Bang's blocks be $b_1,\ldots,b_{n+1}$.
\begin{claim}
    Suppose there exist $A,B \subseteq [n+1]$ with $A \neq B$ and $\varepsilon\geq 0$ such that
$$\sum_{i \in A} b_i=\varepsilon+\sum_{j \in B} b_j.$$Then Xiang Yu can guarantee that $c \leq \frac{1+\varepsilon}{2}$.
\end{claim}
\begin{proof}
    By replacing $A$ and $B$ with $A \setminus B$ and $B \setminus A$ respectively, suppose that $A$ and $B$ are disjoint and not both $\varnothing$. Also suppose that $$\sum_{i \in A} b_i-\sum_{j \in B} b_j$$ is at most $\min_{i \in A} b_i$, else remove the minimal $b_i$ from $A$ to decrease $\varepsilon$. Line up the elements of $A$ end-to-end to form a segment $L_A$ with $|A|-1$ cuts on it, and do the same with $B$ to form $L_B$. Then, along the length of $L_B$, make a cut wherever $L_A$ has a cut. Similarly, along the length of $L_A$, make a cut wherever $L_B$ has a cut, and also make a cut at the end of $L_B$. This is shown below, with black cuts coming from Liu Bang and red cuts coming from Xiang Yu. Overlapping cuts (none shown) are fine, just ignore them.

\medskip
    \begin{figure}[H]
        \centering
        \begin{asy}
            size(8cm);
            draw((0,0)--(10,0), grey);
            draw((0,2)--(12,2), grey);
            draw((2,-0.5)--(2,0.5));
            draw((5,-0.5)--(5,0.5));
            draw((3,1.5)--(3,2.5));
            draw((9,1.5)--(9,2.5));
            draw((2,1.5)--(2,2.5), red);
            draw((5,1.5)--(5,2.5), red);
            draw((3,-0.5)--(3,0.5), red);
            draw((9,-0.5)--(9,0.5), red);
            draw((10,1.5)--(10,2.5), red);
        \end{asy}
    \end{figure}
The result is that Xiang Yu makes at most $|A|+|B|-1$ cuts, and produces $|A|+|B|-1$ pairs of slices of matching length, along with a "leftover" slice $\ell$ of length $\varepsilon$. With the remaining cuts, Xiang Yu cuts each block not indexed in $A$ or $B$ in half. The result is that, aside from $\ell$ all slices can be paired up by length. Then Xiang Yu can employ the following strategy. Before $\ell$ is chosen, he mirrors Liu Bang's moves; whenever Liu Bang chooses a non-$\ell$ slice, Xiang Yu chooses its equal-length pair. At some point, Liu Bang chooses $\ell$; then Xiang Yu switches to greedily choosing slices, which guarantees him at least half of the remaining length. Therefore, Xiang Yu finishes with his total length at most $\varepsilon$ less than Liu Bang's, i.e. $c \leq \frac{1+\varepsilon}{2}$. 
\hfill$\square$
\end{proof}
\noindent
Since there are $2^{n+1}$ subsets of $[n+1]$, it follows that two of them must index $b_\bullet$-sums within $\frac{1}{2^{n+1}-1}$ of each other, so we can apply the claim with $\varepsilon=\frac{1}{2^{n+1}-1}$, which establishes the desired upper bound on $c$. \hfill$\square$

\paragraph{\textcolor{blue}{Strategy for Liu Bang:}} Let's introduce a few more modifications. WLOG force Xiang Yu to also make exactly $n$ cuts. Thus let $s_1,\cdots,s_{2n+1}$ be the slices in decreasing order of length. Abusing notation, also let $s_i$ denote the length of $s_i$; the usage is intuitive and clear in context. Then $c$ equals $\frac{1}{2}$ plus half the quantity
$$S:=(s_1-s_2)+(s_3-s_4)+\cdots+(s_{2n-1}-s_{2n})+s_{2n+1},$$so it suffices to maximize $S$ instead. Scale so that the stick has length $2^{n+1}-1$; now we want to show that the Liu Bang can force $S \geq 1$.

\medskip
Liu Bang's strategy is to cut the stick into blocks of lengths $2^0,\ldots,2^n$. Let $b_i$ be the block of length $2^i$. Construct the following graph $G$ on the $b_i$: for each $j \in [n]$, draw an edge (possibly a self-loop) between the block containing $s_{2j}$ and the block containing $s_{2j-1}$. Then color the block with slice $s_{2n+1}$ red. The number of slices in a block is thus its degree, plus one if it's red. Thus $G$ has $n+1$ vertices and $n$ edges, so there exists a connected component which is a tree. Let $b_k$ be the longest block in this tree. Consider an edge adjacent to $b_k$, formed due to slice $s$ in $b_k$, and slice slice $s'$ in another block, say $b_\ell$. Then we get a contribution of at least $s-s' \leq s-2^{\ell}$ to $S$ (from the $|s-s'|$ term). Summing, we get that the slices in $b_k$ corresponding to edges have total length at most
$$2^0+\cdots+2^{k-1}+C,$$where $C>0$ is an aggregate contribution to $S$. If $C \geq 1$ we are already done. Otherwise, this quantity is strictly less than $2^k$, so we must have at least one more slice. But since the number of slices in $b_k$ is its degree plus $1$ if it's red, we must have exactly one more slice, and it must be $s_{2n+1}$. Furthermore, its length is at least $1-C$. Adding on our earlier contribution (which does not overlap), we get that $S \geq C+(1-C) \geq 1$, as desired. \hfill$\square$

\newpage
\subsubsection{Solution 3}
Let each of the partitioned sticks be of length $a_1, a_2, \dots, a_{2n}$, in decreasing order, as $$a_1 \geq a_2 \geq \dots \geq a_{2n}.$$ Note that then, Liu always takes $a_1, a_3, \dots, a_{2n-1}$. Hence, we must find the minimum of
\[\sum_{1 \leq i = \text{odd} < 2n} a_i\]
\begin{lemma}
    Liu guarantee's maximality when taking sticks in a geometric progression of common ratio $\frac{1}{2^{n+1} -1}$.
\end{lemma}
\begin{proof}
    Liu wants to partition the sticks such that no matter where or how he cuts the larger stick, he can still obtain a net growth. Suppose he creates a cut of two sticks $A < B$. Now, if Xiang leaves $B$ uncut, Liu will grab $B$ first. Otherwise, if he cuts $B$, the worst possible case is that he cuts it into direct halves. To not render $B$ useless, Liu can cut $B$ into double that of $A$, hence it is a net gain greater than or equal to 0. But the thing is, if we have only $A$ and $B$ as our leading sticks, with $B$ double that of $A$, then three new sticks are formed after Xiang cuts, from which by Pigeonhole Principle, Liu obtains the last stick, hence winning the game.

We can guarantee that doubling is the best possible strategy for Liu. If he were to cut triple of quadruple that into pieces $A$ and $B$, Xiang can make some number of cuts with $n$ cuts such that all pieces are useless in comparison to $A$. (Indeed, with the max length 1, there is more than enough $n$ to pass around, as you can check via the base cases.)

\medskip
So, via the doubling strategy, each of the sticks have lengths $\delta(1, 2, 4, \dots, 2^n)$. Here, $\delta$ is our common ratio. This inner sum is $2^{n+1}-1$. Then, $\delta(2^{n+1} -1) = 1$ hence $\delta = \frac{1}{2^{n+1} -1}$, as desired.\hfill$\square$
\end{proof}
\begin{lemma}
    Regardless of Xiang's plays, Liu can guarantee $\frac{2^n}{2^{n+1} -1}$ of minimum length.
\end{lemma}
\begin{proof}
    Suppose he were to make these cuts in the common ratio found in \textcolor{BrightAqua}{\textbf{Lemma 1.3.1}}. Then, given such $n+1$ cuts, the only plausible moveset that Xiang could play is to split all given segments in half. That is, he MUST split the larger $n$ segments in half, whereas the last segment of length 1 is as desired. Then, Liu obtains all odd numbered lengths when ordered $a_1 \rightarrow a_n$ in decreased order, which is $\delta(1+1+2+4+\dots+2^{n-1}) = \delta(2^n - 1 + 1)= \delta \cdot 2^n$, which implies the length $\frac{2^n}{2^{n+1} - 1}$, as desired.\hfill$\square$
\end{proof}
\noindent
We denote $c$ to be the minimum length regardless of Xiang's plays. If we prove that Xiang plays optimally yet we still obtain a minimum length, then we may guarantee that Liu obtains a length greater than or equal to this minimum by definition. That is, through \textcolor{BrightAqua}{\textbf{Lemma 1.3.2}}.
\[c = \boxed{\frac{2^n}{2^{n+1} - 1}} \quad \square\]






\newpage
\section{Solutions to Day II}
\subsection{International Mathematical Olympiad/Problem 4}
\begin{problem*}
    Shan-Yu and Mulan are playing a game. Let $\theta$ be an angle with $0^\circ < \theta < 180^\circ$ known to both players. Initially, Shan-Yu makes a paper triangle $\mathcal{T}$ with measurements of his choice. Then, they repeatedly perform the following steps:
\begin{itemize}
    \item If $\mathcal{T}$ has at least one angle measuring exactly $\theta$, then the game stops and Mulan wins.
    \item Otherwise, Mulan chooses a point $P$ on the perimeter of $\mathcal{T}$, different from its three vertices. She then makes a straight cut from $P$ to the opposite vertex of $\mathcal{T}$, splitting it into two triangles.
    \item Shan-Yu discards one of the two triangles. The remaining triangle becomes the new $\mathcal{T}$.
\end{itemize}
For which real values of $\theta$ can Mulan guarantee her victory in finitely many steps, no matter how Shan-Yu plays?
\end{problem*}
\subsubsection{Solution 1}
The answer is every $\boxed{\theta=180^{\circ}/n}$ for every integers $n\geqslant 2$. We will demonstrate that these work, and the others won't.
\paragraph{\textcolor{blue}{Disprove} $\textcolor{blue}{\boldsymbol{\theta>90^{\circ}:}}$} We say that the triangle is recessive if and only if its largest internal angle is at most $90^{\circ}$.
\begin{claim}
    For any recessive triangle $\mathcal{T}$, we can only split $\mathcal{T}$ into either both recessive triangles or exactly one recessive triangle.
\end{claim}
\begin{proof}
    Suppose that the recessive triangle $\mathcal{T}$ is splitted into $\mathcal{T}_1$ and $\mathcal{T}_2$ by choosing a striaght cut from the vertex $A$ to the point $P$ on the perimeter. Suppose that the other two vertices are $B$ and $C$. If $\angle{APB}=90^{\circ}$, then both $\mathcal{T}_1$ and $\mathcal{T}_2$ are recessive, and if $\angle{APB}<90^{\circ}$, then $\angle{APC}>90^{\circ}$; which makes exactly one of $\mathcal{T}_1$ or $\mathcal{T}_2$ is recessive. So, the claim is completed. Suppose that the recessive triangle $\mathcal{T}$ is splitted into $\mathcal{T}_1$ and $\mathcal{T}_2$ by choosing a striaght cut from the vertex $A$ to the point $P$ on the perimeter. Suppose that the other two vertices are $B$ and $C$. If $\angle{APB}=90^{\circ}$, then both $\mathcal{T}_1$ and $\mathcal{T}_2$ are recessive, and if $\angle{APB}<90^{\circ}$, then $\angle{APC}>90^{\circ}$; which makes exactly one of $\mathcal{T}_1$ or $\mathcal{T}_2$ is recessive. So, the claim is completed.
    
    \hfill$\square$
\end{proof}
\noindent
Therefore, Shan-Yu can start with the equilateral triangle (which is recessive). By the above claim, notice that regardless of how Mulan splits, Shan-Yu can guaranteed to discard the triangle $\mathcal{T}_1$ so that the remaining triangle is always recessive. Hence, Mulan can not have $\theta>90^{\circ}$ as one of the internal angles. This makes every $\theta>90^{\circ}$ fail to satisfy the problem's condition.
\paragraph{\textcolor{blue}{Prove every} $\textcolor{blue}{\boldsymbol{\theta=180^{\circ}/n}:}$}
It is clear that $\theta=90^{\circ}$ works as Mulan can split the triangle but making a straight cut from the vertex having the largest angle to its foot of altitude. Hence, we may assume that $\theta<90^{\circ}$.
\begin{claim}
    Suppose that a triangle $\mathcal{T}$ has the smallest angle $\angle{A}<\theta$, let $B\neq A$ be the vertex such that $\angle{B}<90^{\circ}$, and $C\neq A,B$ be another vertex, then we can choose such point $P$ on $AB$ satisfying $\angle{CPB}=\theta$.
\end{claim}
\begin{proof}
    Notice that $\angle{A}<180^{\circ}-\angle{B}$ because $\angle{A}+\angle{B}<180^{\circ}$, so, $\angle{CPB}$ can take every value in $(\angle{A},180^{\circ}-\angle{B})$ by animating $P\in AB$. The condition $\angle{A}<\theta$ and $\theta<180^{\circ}-\angle{B}$ are true as $\theta+\angle{B}<90^{\circ}+90^{\circ}=180^{\circ}$; that is $\theta\in(\angle{A},180^\circ{-\angle{B}})$. Now we get the desired result.\hfill$\square$
\end{proof}
\begin{figure}[H]
    \centering
    \begin{asy}
        size(8cm);

        import geometry;
        import graph;

        pen pen_ff0000 = RGB(255, 0, 0);
        pen pen_0000ff = RGB(0, 0, 255);

        path anglemark(pair A, pair B, pair C, real t=8 ... real[] s)
        {
        real markscalefactor=0.07;
        pair M,N,P[],Q[];
        path mark;
        int n=s.length;
        M=t*markscalefactor*unit(A-B)+B;
        N=t*markscalefactor*unit(C-B)+B;
        for (int i=0; i<n; ++i)
        {
        P[i]=s[i]*markscalefactor*unit(A-B)+B;
        Q[i]=s[i]*markscalefactor*unit(C-B)+B;
        }
        mark=arc(B,M,N);
        for (int i=0; i<n; ++i)
        {
        if (i%2==0)
        {
        mark=mark--reverse(arc(B,P[i],Q[i]));
        }
        else
        {
        mark=mark--arc(B,P[i],Q[i]);
        }
        }
        if (n%2==0 && n!=0)
        mark=(mark--B--P[n-1]);
        else if (n!=0)
        mark=(mark--B--Q[n-1]);
        else mark=(mark--B--cycle);
        return mark;
        }

        pair A = (1.75,0);
        pair B = (-1,0);
        pair C = (-0.3,1.767);
        pair P = (1.3,0);

        dot("$A$",A,SE);
        dot("$B$",B,SW);
        dot("$C$",C,N);
        dot("$P$",P,S);

        draw(A--B--C--cycle,pen_ff0000);
        draw(C--P,pen_0000ff);

        label("$\theta$", anglemark(C,P,B,0.01), 3.4W+1.5N, fontsize(9pt));
    \end{asy}
\end{figure}
\noindent
First, notice that Mulan can repeatedly bisects some angle of $\mathcal{T}$ so that the resulting $\mathcal{T}$ has the smallest internal smaller than $180^{\circ}/n$. Suppose that $\mathcal{T}$ has the smallest $\angle{A}<180^{\circ}/n$ and $B\neq A$ be one vertex of $\mathcal{T}$ such that $\angle{B}<90^{\circ}$ (which exists as we can't have two vertices that their angles measure at least $90^{\circ}$), and let $C\neq A,B$ be another vertex of $\mathcal{T}$. From the above claim, we can pick the point $P$ on the line $AB$ such that $\angle{CPB}=180^{\circ}/n$ (and $\angle{CPA}=(n-1)180^{\circ}/n$.)

\medskip
Hence, each step Shan-Yu is forced to remove the triangle with $180^{\circ}/n$ internal angle, and then Mulan can split the angle $(n-1)180^{\circ}/n$ to $180^{\circ}/n$ and $(n-2)180^{\circ}/n$, and so on. Finally, one of the internal angle eventually becomes $180^{\circ}/n$. Therefore, Mulan can guarantee to win with these angles.

\medskip
Here is the demonstration on how Mulan plays with $5\theta=180^{\circ}$ (take note that this is what we got after we spam the angle bisection from our previous claim.)

\medskip
\begin{figure}[H]
    \centering
    \begin{asy}
        size(16cm);
        import geometry;
        import graph;

        path anglemark(pair A, pair B, pair C, real t=8 ... real[] s)
        {
        real markscalefactor=0.07;
        pair M,N,P[],Q[];
        path mark;
        int n=s.length;
        M=t*markscalefactor*unit(A-B)+B;
        N=t*markscalefactor*unit(C-B)+B;
        for (int i=0; i<n; ++i)
        {
        P[i]=s[i]*markscalefactor*unit(A-B)+B;
        Q[i]=s[i]*markscalefactor*unit(C-B)+B;
        }
        mark=arc(B,M,N);
        for (int i=0; i<n; ++i)
        {
        if (i%2==0)
        {
        mark=mark--reverse(arc(B,P[i],Q[i]));
        }
        else
        {
        mark=mark--arc(B,P[i],Q[i]);
        }
        }
        if (n%2==0 && n!=0)
        mark=(mark--B--P[n-1]);
        else if (n!=0)
        mark=(mark--B--Q[n-1]);
        else mark=(mark--B--cycle);
        return mark;
        }

        picture[] pics = {};
        int n=5;
        for(int i=0; i<n; ++i){
        picture pic;
        pics.push(pic);
        }

        pen pen_ff0000 = RGB(255, 0, 0);
        pen pen_ff6699 = RGB(255, 102, 153);
        pen lg = RGB(180, 180, 180);

        pair point_A = (-0.4460910116830806, 0.5662292276968843);
        pair point_B = (-0.7500000000000000, 0.0000000000000000);
        pair point_C = (1.5000000000000000, 0.0000000000000000);
        pair point_P_1 = (0.3325800288679810, 0.0000000000000000);
        pair point_P_2 = (0.6662226038593585, 0.2425935520248511);
        pair point_P_3 = (0.4333932972028370, 0.3103369194736868);
        pair point_P_4 = (0.2107721242930510, 0.3751101567215501);
        path ab = point_A--point_B;
        path bc = point_B--point_C;
        path ca = point_C--point_A;
        path ap1 = point_A--point_P_1;
        path bp1 = point_B--point_P_1;
        path p2c = point_P_2--point_C;
        path cp1 = point_C--point_P_1;
        path p1p3 = point_P_1--point_P_3;
        path p1p2 = point_P_1--point_P_2;
        path p2p3 = point_P_2--point_P_3;
        path p3a = point_P_3--point_A;
        path p4p1 = point_P_4--point_P_1;
        path ap2 = point_A--point_P_2;

        // fig 0
        fill(pics[0],point_A--point_B--point_C--cycle, palered+white);

        draw(pics[0],ab, pen_ff0000);
        draw(pics[0],bc, pen_ff0000);
        draw(pics[0],ca, pen_ff0000);
        draw(pics[0],ap1, pen_ff6699+dashed);

        dot(pics[0],point_A);
        dot(pics[0],point_B);
        dot(pics[0],point_C);
        dot(pics[0],point_P_1);

        label(pics[0],"$C$", point_A, dir(point_A));
        label(pics[0],"$B$", point_B, dir(point_B));
        label(pics[0],"$A$", point_C, dir(point_C));
        label(pics[0],"$P_{1}$", point_P_1, S);

        // fig 1
        fill(pics[1],point_A--point_P_1--point_C--cycle, palered+white);

        draw(pics[1],ab, lg);
        draw(pics[1],bp1, lg);
        draw(pics[1],ap1, pen_ff0000);
        draw(pics[1],ca, pen_ff0000);
        draw(pics[1],cp1, pen_ff0000);
        draw(pics[1],p1p2,pen_ff6699+dashed);

        dot(pics[1],point_A);
        dot(pics[1],point_B,lg);
        dot(pics[1],point_C);
        dot(pics[1],point_P_1);
        dot(pics[1],point_P_2);

        label(pics[1],"$C$", point_A, dir(point_A));
        label(pics[1],"$B$", point_B, dir(point_B), lg);
        label(pics[1],"$A$", point_C, dir(point_C));
        label(pics[1],"$P_{1}$", point_P_1, S);
        label(pics[1],"$P_2$", point_P_2,NE);

        // fig 2
        fill(pics[2],point_A--point_P_1--point_P_2--cycle, palered+white);

        draw(pics[2],ab, lg);
        draw(pics[2],bp1, lg);
        draw(pics[2],ap1, pen_ff0000);
        draw(pics[2],p2c, lg);
        draw(pics[2],cp1, lg);
        draw(pics[2],p1p2,pen_ff0000);
        draw(pics[2],ap2,pen_ff0000);
        draw(pics[2],p1p3,pen_ff6699+dashed);

        dot(pics[2],point_A);
        dot(pics[2],point_B,lg);
        dot(pics[2],point_C,lg);
        dot(pics[2],point_P_1);
        dot(pics[2],point_P_2);
        dot(pics[2],point_P_3);

        label(pics[2],"$C$", point_A, dir(point_A));
        label(pics[2],"$B$", point_B, dir(point_B), lg);
        label(pics[2],"$A$", point_C, dir(point_C), lg);
        label(pics[2],"$P_{1}$", point_P_1, S);
        label(pics[2],"$P_2$", point_P_2,NE);
        label(pics[2],"$P_3$", point_P_3,N);

        // fig 3
        fill(pics[3],point_A--point_P_1--point_P_3--cycle, palered+white);

        draw(pics[3],ab, lg);
        draw(pics[3],bp1, lg);
        draw(pics[3],ap1, pen_ff0000);
        draw(pics[3],p2c, lg);
        draw(pics[3],cp1, lg);
        draw(pics[3],p1p2,lg);
        draw(pics[3],p1p3,pen_ff0000);
        draw(pics[3],p3a,pen_ff0000);
        draw(pics[3],p4p1,pen_ff6699+dashed);
        draw(pics[3],p2p3,lg);

        dot(pics[3],point_A);
        dot(pics[3],point_B,lg);
        dot(pics[3],point_C,lg);
        dot(pics[3],point_P_1);
        dot(pics[3],point_P_2,lg);
        dot(pics[3],point_P_3);
        dot(pics[3],point_P_4);

        label(pics[3],"$C$", point_A, dir(point_A));
        label(pics[3],"$B$", point_B, dir(point_B), lg);
        label(pics[3],"$A$", point_C, dir(point_C), lg);
        label(pics[3],"$P_{1}$", point_P_1, S);
        label(pics[3],"$P_2$", point_P_2,NE,lg);
        label(pics[3],"$P_3$", point_P_3,N+0.3E);
        label(pics[3],"$P_4$", point_P_4,N);

        // angle label
        label(pics[0], "$\theta$", anglemark(point_A,point_P_1,point_B), 0.7NW, fontsize(9pt));
        label(pics[1], "$\theta$", anglemark(point_C,point_P_1,point_P_2), 0.3S+0.6E, fontsize(9pt));
        label(pics[2], "$\theta$", anglemark(point_P_3,point_P_1,point_P_2,0.0001), 1.4E +1.8N, fontsize(9pt));
        label(pics[3], "$\theta$", anglemark(point_A,point_P_1,point_P_4,0.0001), 3N+0.2E, fontsize(9pt));

        // add
        add(pics[0].fit(6cm,6cm), (0,0));
        add(pics[1].fit(6cm,6cm), (6,0));
        add(pics[2].fit(6cm,6cm), (0,-2));
        add(pics[3].fit(6cm,6cm), (6,-2));
    \end{asy}
\end{figure}
\paragraph{\textcolor{blue}{Disprove every} $\textcolor{blue}{\boldsymbol{\theta\neq 180^{\circ}/n}:}$}
We say that an angle is good if and only if its measure is of the form $k\cdot \theta$ for some positive integer $k$. We start by proving the following claim:
\begin{claim}
    Suppose that $\mathcal{T}$ contains no good internal angles, then for any splits, there are at most one triangle that could contain some good angle.
\end{claim}
\begin{proof}
    Suppose that $A,B$ and $C$ are vertices of $\mathcal{T}$ and the cut splits $\mathcal{T}$ into $\triangle{ABP}$ and $\triangle{ACP}$. Consider the following cases:
    \begin{itemize}
        \item If $\angle{BAP}$ is good, then $\angle{APC}$ is not good and $\angle{PAC}$ is not good. So, $\triangle{APC}$ can not have any good angles.
        \item If $\angle{APB}$ is good, then $\angle{APC}$ is not good and $\angle{APC}$ is not good. So, $\triangle{APC}$ can not have any good angles.
    \end{itemize}
    Similarly, if at least one of $\angle{PAC}$ or $\angle{APC}$ is good, then $\triangle{APB}$ can not contain any good angles. \hfill$\square$
\end{proof}
\noindent
Shan-Yu may start with the equilateral triangle so that none of the internal angles of $\mathcal{T}$ is good. By using the claim above, after each splits, Shan-Yu can guarantee to remove one triangle so that the new $\mathcal{T}$ does not contain any good internal angles, and this process can be done indefinitely. So, Mulan can not guarantee her win in finitely many steps.
\newpage
\subsection{International Mathematical Olympiad/Problem 5}
\begin{problem*}
    Let $\mathbb{R}_{>0}$ be the set of positive real numbers. Determine all functions $f \colon \mathbb{R}_{>0} \to \mathbb{R}_{>0}$ such that
\[
\sqrt{\frac{x^2 + f(y)^2}{2}} \ge \frac{f(x) + y}{2} \ge \sqrt{x f(y)}
\]
for every $x, y \in \mathbb{R}_{>0}$.
\end{problem*}
\subsubsection{Solution 1}
The answer is $\boxed{f(x)=x+c}$, which works because of QM-AM-GM inequality.

\medskip
\noindent
First, plugging in $x=f(y)$ gives that
$$f(y) \geq \frac{f(f(y)) + y}2 \geq f(y)$$so $\boxed{f(f(y)) = 2f(y)-y}$ for all $y$. Iterating gives
$$f^n(y) = y + n(f(y)-y)$$for all $y$, and sending $n\to\infty$ gives that $\boxed{f(y)\geq y}$ for all $y$.

\medskip
\noindent
Here is the main claim.
\begin{claim}
     $f\circ f$ is non-decreasing. In particular, if $x\geq y$, then $f(f(x))\geq f(f(y))$.
\end{claim}
\begin{proof}
    Replacing $x$ with $f(x)$ in the left inequality gives
\begin{align*}
\sqrt{\frac{f(x)^2 + f(y)^2}2} &\geq \frac{f(f(x)) + y}2 = \frac{2f(x)-x+y}2 \\
2f(x)^2 + 2f(y)^2 &\geq (2f(x)-x+y)^2.
\end{align*}Similarly, $2f(x)^2 + 2f(y)^2 \geq (2f(y)-y+x)^2$. Summing gives
\begin{align*}
4f(x)^2 + 4f(y)^2 &\geq (2f(x)-x+y)^2 + (2f(y)-y+x)^2 \\
&\geq 4f(x)^2 + 4f(y)^2 - 4(f(x)-f(y))(x-y) + 2(x-y)^2,
\end{align*}which rearranges to
$$2(x-y)(2f(x)-2f(y)-x+y) \geq 0.$$Thus, $2f(x)-x \geq 2f(y)-y$, which simplifies to $f(f(x))\geq f(f(y))$. \hfill$\square$
\end{proof}
\begin{claim}
     $f(x)-x$ is non-decreasing.
\end{claim}
\begin{proof}
    From above, we have that $f^{2n}$ is non-decreasing. Thus, for all $x\geq y$, we have
$$x+2n(f(x)-x) \geq y+2n(f(y)-y),$$which implies that $f(x)-x\geq f(y)-y$ after sending $n\to\infty$. \hfill$\square$
\end{proof}
\begin{claim}
     $f(x)-x$ is either $0$ or $a$ for some constant $a$.
\end{claim}
\begin{proof}
    Assume that there exists $x<y$ such that $f(x)=x+a$ and $f(y)=y+b$ where $a,b>0$. Using $f(f(x)) = 2f(x)-x$, we get $f(x+a)=x+2a$, and by induction, $f(x+na) = x+(n+1)a$. Thus, $t\mapsto f(t)-t$ is constant in interval $[x+na, x+(n+1)a]$ for all $a$ (since we already knew it's monotone), which implies that $f(y)-y = f(x)-x$, a contradiction. \hfill$\square$
\end{proof}
\noindent
Now, we only have to deal with pointwise trap. If $f(x)=x$ and $f(y)=y+a$, then
$$\frac{x^2+y^2}{2} \geq \left(\frac{x+y+a}2\right)^2 \geq \frac{(x+y)^2+a^2}{4},$$which implies that that $x-y > 0.01a$. In particular, $f(x)-x$ is constant in any interval of length $0.01a$, so by induction, $f(x)-x$ is constant.

\newpage
\subsubsection{Solution 2}
The solutions are of the form $\boxed{f(x) = x + c}$, and these clearly work by power means inequality.

\medskip
Firstly, let's try to get an equality out of the inequalities. By looking at $P(f(x), x)$, the QM and the GM are the same(they're both $f(x)$), and this implies
\[ \frac{f(f(x)) + x}{2} = f(x) , \]and in particular this means the sequence $x, f(x), f(f(x)), \dots$ is an arithmetic progression. Let $c_x$ denote the common difference of this arithmetic progression for $x$. $c_x$ must be nonnegative since otherwise some $f^n(x)$ would have to be negative, contradiction.
\begin{claim}
     $f$ is increasing.
\end{claim}
\begin{proof}
    By plugging in every combination of $x$, $y$, $f(x)$ and $f(y)$ until something works, we see that if $ x > y$, $P(x, y)$ shows
\begin{align*}
\frac{f(f(x)) + y}{2} \ge \sqrt{f(x)f(y)} \\
&\implies \frac{2f(x) - x + y}{2} \ge \sqrt{f(x)f(y)} \\
&\implies \frac{2f(x)}{2} \ge \sqrt{f(x)f(y)} \\ 
&\implies f(x) > f(y).
\end{align*}Thus $f$ is increasing. \hfill$\square$
\end{proof}
\begin{claim}
    There cannot be two distinct positive values of the common difference.
\end{claim}
\begin{proof}
    Assume $c_x > c_y$, for some values of $x$ and $y$. Our aim will be to find some $r \coloneq x + ic_x$ and $s \coloneq y + jc_y$ such that $r < s$ but $f(r) > f(s)$. Since $r$ and $s$ are part of the arithmetic progression, $f(r) > f(s)$ means $r + c_x > s + c_y$, so it suffices to show that we can choose $r$ and $s$ to be within $c_x - c_y$ of each other. Rewriting this, it becomes $(j - 1)c_y - (i - 1)c_x > x - y$, which is possible by just making $j$ as large as we can. \hfill$\square$
\end{proof}
\noindent
Thus, we have $f(x) = x$ or $f(x) = x + c$ for some global choice of $c$. To deal with the pointwise trap , assume $f(x) = x$ but $f(y) = y + c$.
Then $P(x, y)$ reads
\[ \frac{x + y}{2} \ge 4(x(y + c))  \Leftrightarrow (x - y)^2 \ge 4xc , \]but we can take $x$ and $y$ to be arbitrarily close to get a contradiction here.

\newpage
\subsubsection{Solution 3}
Taking $x = f(y)$ yields $f(f(y)) + y = 2f(y)$. By taking $x = f(y) + \varepsilon$ we aim to find the derivative of $f$ over points in the image of $f$. For this notice that:
$$\sqrt{\frac{(f(y)+\varepsilon)^2+f(y)^2}{2}} \geq \frac{f(f(y) + \varepsilon)+y}{2} \geq \sqrt{(f(y)+\varepsilon) f(y)}$$$$\sqrt{f(y)^2 + \varepsilon f(y) + \frac{\varepsilon^2}{2}} - f(y) \geq \frac{f(f(y) + \varepsilon)+y - 2f(y)}{2} \geq \sqrt{f(y)^2+\varepsilon f(y)}-f(y)$$Dividing by $\varepsilon$, taking $\varepsilon \to 0$ and because $y - 2f(y) = -f(f(y))$ we get:
$$1 \geq \lim_{\varepsilon \to 0} \frac{f(f(y) + \varepsilon) - f(f(y))}{\varepsilon} \geq 1$$And $f(f(y)+\varepsilon) = f(f(y)) + \varepsilon + o_y(\varepsilon) \varepsilon$ where $o_y(\delta) \to 0$ as $\delta \to 0$. Now using the identity from the start:
$$f(f(y+ \varepsilon)) + y + \varepsilon = 2f(y + \varepsilon)$$And calling $d_{\varepsilon,y} = f(y+\varepsilon) - f(y)$ we have:
$$f(f(y)) + d_{\varepsilon,y} + o_y(d_{\varepsilon,y})d_{\varepsilon,y} + y + \varepsilon = 2f(y + \varepsilon)$$$$d_{\varepsilon,y} + o_y(d_{\varepsilon,y})d_{\varepsilon,y} + \varepsilon = 2 f(y + \varepsilon) - y - f(f(y)) = 2 d_{\varepsilon,y}$$$$\varepsilon = d_{\varepsilon,y} (1 - o_y(d_{\varepsilon,y}))$$But notice that, from before, if $d>0$ we can bound the incremental cocient by:
$$  o_y(d) \geq \frac{2}{d} (\sqrt{f(y)^2+d f(y)} - f(y) - \frac{d}{2} ) = - \frac{ \frac{d}{2}}{\sqrt{f(y)^2+d f(y)} + f(y) + \frac{d}{2}} \geq - \frac{\frac{d}{2}}{f(y) + \frac{d}{2}}$$And:
$$o_y(d) \leq \frac{2}{d}(\sqrt{f(y)^2 + d f(y) + \frac{d^2}{2}} - f(y)  - \frac{d}{2}) =  \frac{\frac{d}{2}}{\sqrt{f(y)^2 + d f(y) + \frac{d^2}{2}} + f(y) + \frac{d}{2}} \leq \frac{\frac{d}{2}}{f(y) + \frac{d}{2}}$$So:
$$|1 - o_y(d_{\varepsilon, y})| \geq 1 - |o_y(d_{\varepsilon, y})| \geq \frac{f(y)}{f(y) + \frac{d_{\varepsilon, y}}{2}}$$Now, taking $\varepsilon \to 0$ in:
$$\varepsilon = d_{\varepsilon,y} |1 - o_y(d_{\varepsilon,y})| \geq  d_{\varepsilon,y} \frac{f(y)}{f(y) + \frac{d_{\varepsilon,y}}{2}} \Rightarrow d_{\varepsilon,y} \leq \frac{\varepsilon f(y)}{f(y) - \frac{\varepsilon}{2}}$$Yields $d_{\varepsilon,y} \to 0$. With the case of $d<0$ being almost analogous, we have then that:
$$1 = \lim_{\varepsilon \to 0}  \frac{d_{\varepsilon,y}}{ \varepsilon} (1 - o_y(d_{\varepsilon,y})) = f'(y)$$At every point $ y \in \mathbb{R}_{>0}$, so the only candidates are $f(y) = y+c$ for some $c \geq 0$ which can be checked to be solutions by simple HM-AM-GM inequality.
\newpage
\subsection{International Mathematical Olympiad/Problem 6}
\begin{problem*}
    Let $a_1, a_2, a_3, \dots$ be an infinite sequence of positive integers greater than $1$. Suppose that for all positive integers $n$, the number $a_{n+1}$ is the smallest positive integer greater than $a_n$ such that $\gcd(a_{n+1}, a_i) > 1$ for every $i = 1, 2, \dots, n$. Prove that there exist positive integers $T$ and $L$ such that
\[
a_{n+T} = a_n + L
\]
for every positive integer $n$.
\end{problem*}
\subsubsection{Solution 1}
Let $P(m)$ be the set of all prime divisors of $m$ and consider
\[ \mathcal{S} := \{ m \in \mathbb{Z}_{>0} :  \gcd(m, a_i) > 1, \quad \forall i \ge 1\}.\]
\begin{claim}
    $a_1, a_2, \dots$ are precisely the increasing enumeration of the elements of $\mathcal{S}$ which are at least $a_1$.
\end{claim}
\begin{proof}
    Note that if $i < j$, the choice of $a_j$ forces $\gcd(a_j, a_i) > 1$, which means that $a_j \in \mathcal{S}$ for any $j$. Now, suppose that there exists $m \in \mathcal{S}$ which is not enumerated, then there exists some $n$ for which $a_n < m < a_{n + 1}$. However, this means that $\gcd(m, a_i) > 1$ for all $1 \le i \le n$, since $m \in \mathcal{S}$, while $m < a_{n + 1}$, which contradicts the minimality of $a_{n + 1}$.\hfill$\square$
\end{proof}
\noindent
We can now phrase this condition above in terms of combinatorics of set.

\medskip
\noindent
Let $\mathcal{F}$ denote the family of all finite sets $S$ of primes for which $S \cap P(a_n) \not= \varnothing$ for all $n \ge 1$.
\begin{claim}
    $\mathcal{F}$ is intersecting.
\end{claim}
\begin{proof}
    Suppose otherwise, that there exists $X, Y\in \mathcal{F}$ for which $X \cap Y = \varnothing$. Note that neither $X$ and $Y$ can be empty by definition. Therefore, we may choose $a$ and $b$ large enough such that
\[ x := \left( \prod_{p \in X} p \right)^{a} \ \text{and} \ y := \left( \prod_{p \in Y} p \right)^{b} \]while $x, y > a_1$. We see that $P(x) = X$ and $P(y) = Y$. Since $X, Y \in \mathcal{F}$, we see that $x, y \in \mathcal{S}$. Since $x, y > a_1$, this means that $x, y$ appears in the enumeration of the sequence $a_1, a_2, \dots$. However, $\gcd(x, y) = 1$ since $X \cap Y = \varnothing$, which is a contradiction.\hfill$\square$
\end{proof}
\noindent
Consider the minimal members of $\mathcal{F}$, i.e., the elements in $\mathcal{F}$ for which no proper subset of it belongs to $\mathcal{F}$.

\medskip
\noindent
The crux of this entire problem lies on the following main claim.
\begin{claim}[Main Claim]
    If a prime $p$ belongs to a minimal member of $\mathcal{F}$, then $p \le a_1^2$.
\end{claim}
\begin{proof}
    Take any such prime $p$. Among all minimal member of $\mathcal{F}$ containing $p$, choose one for which the product of its elements is the smallest, call this set $M$. If $M = \{ p \}$, then since $M \cap P(a_1) \not= \varnothing$, we must have $p \mid a_1$ and thus $p \le a_1 \le a_1^2$. Therefore, we may assume that $M \setminus \{ p \} \not= \varnothing$, and set $M' := M \setminus \{ p \}$ with $h := \prod_{q \in M'} q$. Since $M$ is minimal, $M' \notin \mathcal{F}$.
    \begin{itemize}
        \item If $h > a_1$, we see that since $M' \notin \mathcal{F}$, we must have $h \notin \mathcal{S}$. This means that there exists $i$ such that $a_i < h$ and $\gcd(a_i, h) = 1$, i.e., $P(a_i) \cap H = \varnothing$. However, $P(a_i) \cap M \not= \varnothing$ and thus $p \in P(a_i)$. Choose a minimal member $N$ of $\mathcal{F}$ such that $N \subseteq P(a_i)$. As $N, M \in \mathcal{F}$, $N \cap M \not= \varnothing$. However, $N \subseteq P(a_i)$ and $P(a_i) \cap M' = \varnothing$, and so $N \cap M' = \varnothing$. This means that $p \in N$. However, the product of every element in $N$ is at most $a_i < h$, due to $N \subseteq P(a_i)$. This gives us a contradiction of our choice $M$.
        \item If $h \le a_1$, let $k$ be the smallest positive integer such that $x := h^k > a_1$. By minimality of $k$, we see that $x \le a_1^2$. We see that $P(x) = M' \notin \mathcal{F}$ and so $x \notin \mathcal{S}$. This means that there exists $i$ such that $a_i < x$ and $\gcd(a_i, x) = 1$, i.e., $P(a_i) \cap M' = \varnothing$. Since $P(a_i) \cap M \not= \varnothing$, we know that $p \in P(a_i)$, which means that $p \le a_i < x \le a_1^2$.\hfill$\square$
    \end{itemize}
\end{proof}
\noindent
Now, let $\mathcal{P} := \{ p \ \text{is prime and } p \le a_1^2 \}$. By our previous claim, every minimal member of $\mathcal{F}$ is a subset of $\mathcal{P}$, and thus there are only finitely many minimal members. Suppose these are $M_1, M_2, \dots, M_r$. For each $1 \le j \le r$, we define
\[ d_j := \prod_{p \in M_j} p. \]
\begin{claim}
    We have
\[ \mathcal{S} = \bigcup_{j= 1}^r \{ m \in \mathbb{Z}_{>0} : d_j \mid m \} \]
\end{claim}
\begin{proof}
    Indeed, we note that
\begin{align*}
n \in \mathcal{S} &\iff P(n) \in \mathcal{F} \\ 
&\iff P(n) \ \text{contain some minimal member } M_j \\ 
&\iff d_j \mid n \ \text{for some } j.
\end{align*}\hfill$\square$
\end{proof}
\noindent
To finish, let $L := \operatorname{lcm}(d_1, \dots, d_r)$. We see that $m \in \mathcal{S} \iff m + L \in \mathcal{S}$ for any $m \ge 1$. This means that membership in $\mathcal{S}$ is periodic modulo $L$. Define $T := |\mathcal{S} \cap [L]|$. Since every interval of $L$ consecutive integers contains exactly $T$ elements of $\mathcal{S}$, and $a_n \in \mathcal{S}$, we see that $a_n + L \in \mathcal{S}$. Furthermore, we see that the sequence $a_n$ is just the increasing enumeration of the elements of $\mathcal{S}$ above $a_1$. As there are precisely $T$ elements of $\mathcal{S}$ lie on $(a_n, a_n + L]$, the $T$-th term after $a_n$ is precisely $a_n + L$, i.e., $a_{n + T} = a_n + L$, for all $n \ge 1$, which is precisely what we wanted.
\newpage
\subsubsection{Solution 2}
Let us define an integer $x$ to be $n$-valid if $\gcd(x, a_i) > 1$ for all $1 \leqslant i \leqslant n$. By definition, $$a_{n+1} = \min \{ x \in \mathbb{Z} \mid x > a_n \text{ and } x \text{ is } n\text{-valid} \}.$$
\begin{claim}
    The sequence of differences $a_{n+1} - a_n$ is strictly bounded.
\end{claim}
\begin{proof}
    Let $P_1$ be the set of distinct prime factors of $a_1$. We define the radical of $a_1$ as $G = \prod_{p \in P_1} p$. Because the condition $\gcd(a_i, a_1) > 1$ must hold for all $i \geqslant 1$, every term $a_i$ must possess at least one prime factor $q \in P_1$. Consequently, $q \mid G$.

\medskip
Consider any multiple of $G$, say $mG$ for some $m \in \mathbb{Z}^+$. For any $i \geqslant 1$, since $a_i$ is divisible by some $q \in P_1$ and $q \mid G$, it follows that $q \mid mG$. Hence, $\gcd(mG, a_i) \geqslant q > 1$. This implies that every multiple of $G$ is $n$-valid for all $n$.

\medskip
At step $n$, there exists at least one multiple of $G$ in the interval $(a_n, a_n + G]$. Because $a_{n+1}$ is the smallest $n$-valid integer strictly greater than $a_n$, we must have:
$$ a_{n+1} \leqslant a_n + G \implies a_{n+1} - a_n \leqslant G $$Thus, the gaps between consecutive terms are bounded by $G$. \hfill$\square$
\end{proof}
\begin{claim}
    Let $\mathcal{P}$ denote the set of all primes that divide at least one term in the sequence $(a_n)_{n=1}^\infty$. The set $\mathcal{P}$ is finite.
\end{claim}
\begin{proof}
    Assume for the sake of contradiction that $\mathcal{P}$ is infinite. Then there exist arbitrarily large primes in $\mathcal{P}$. Let $p \in \mathcal{P}$ be a prime strictly greater than $G$.

\medskip
Let $a_n$ be the very first term in the sequence that is divisible by $p$. This dictates that for all $i < n$, $p \nmid a_i$. Thus, the prime $p$ does not contribute to the condition $\gcd(a_n, a_i) > 1$ for any $i < n$.

\medskip
We can write $a_n = p \cdot k$ for some positive integer $k$. Because $p$ is extraneous to the $n$-validity of $a_n$ with respect to the previous terms, the integer $k$ itself must share a prime factor with every $a_i$ for $i < n$. Ergo, $k$ is also $(n-1)$-valid.

\medskip
Any multiple of an $(n-1)$-valid integer is evidently $(n-1)$-valid. Therefore, the integers $k, 2k, 3k, \dots, p k = a_n$ are all $(n-1)$-valid. Since $a_n$ is defined as the strictly smallest $(n-1)$-valid integer greater than $a_{n-1}$, the preceding multiple $(p-1)k$ must be less than or equal to $a_{n-1}$. This yields:
$$ a_n - a_{n-1} \leqslant p k - (p-1)k = k $$However, we established in Claim 1 that $a_n - a_{n-1} \leqslant G$, and since $k \leqslant a_n - a_{n-1}$, we deduce $k \leqslant G$.

\medskip
Because $k \leqslant G$, $k$ is bounded. Since $\mathcal{P}$ is infinite, we can select our initial prime $p$ to be arbitrarily large, specifically $p > G^2$. However, $a_n = p \cdot k$, and for this large $p$, we must consider whether $p$ ever becomes an essential prime factor for future terms. For $p$ to be strictly required by some future term $a_m$ ($\displaystyle m > n$), $a_m$ must rely on $p$ to share a factor with $a_n$. This implies $a_m$ must be a multiple of $p$.

\medskip
The next available multiple of $p$ is $p(k+1) = a_n + p$. Since $p > G^2 \geqslant G$, we have $a_n + p > a_n + G$. However, the interval $(a_n, a_n + G]$ is guaranteed to contain a multiple of $G$, which is valid for all time. Thus, the sequence will inevitably pick the multiple of $G$ (or an even smaller valid number) long before it ever reaches the next multiple of $p$. Consequently, the prime $p$ is entirely redundant for all future terms.

\medskip
If $p$ is redundant for all $\displaystyle m > n$, and $p$ was redundant for $a_n$ itself (as $k$ was already sufficient), the minimality of $a_n$ is directly contradicted unless $p \leqslant G$. Hence, no such arbitrarily large prime $p$ can exist, and $\mathcal{P}$ must be finite. \hfill$\square$
\end{proof}
\begin{theorem}
    There exist positive integers $T$ and $L$ such that $a_{n+T} = a_n + L$ for every positive integer $n$.
\end{theorem}
\begin{proof}
    By \textcolor{GlacierBlue}{\textbf{Claim 2.3.6}}, the set of primes $\mathcal{P}$ utilized by the entire sequence is finite. Let $$M = \prod_{p \in \mathcal{P}} p.$$

\medskip
The condition that an integer $x$ shares a prime factor with $a_i$ is completely determined by the prime factors of $x$. Since we solely concern ourselves with primes in $\mathcal{P}$, the statement $\gcd(x, a_i) > 1$ is unequivocally equivalent to $\gcd(x, a_i, M) > 1$.

\medskip
Consequently, the $(n-1)$-validity of any candidate integer $x$ depends entirely upon its congruence class modulo $M$. Let $\mathcal{S}_n$ be the set of residue classes modulo $M$ that satisfy the validity condition at step $n$. As $n$ increases, the condition becomes more stringent, meaning $\mathcal{S}_{n+1} \subseteq \mathcal{S}_n$.

\medskip
Because there are only finitely many subsets of residue classes modulo $M$, the sequence of sets $\mathcal{S}_n$ must eventually stabilize. That is, there exists some index $N$ such that for all $n \geqslant N$, $\mathcal{S}_n = \mathcal{S}^*$.

\medskip
For $n \geqslant N$, $a_{n+1}$ is simply the smallest integer strictly greater than $a_n$ whose congruence class modulo $M$ belongs to the fixed, non-empty set $\mathcal{S}^*$. This dictates that for $n \geqslant N$, the sequence $a_n$ is formed by enumerating the elements of a periodic set of integers.

\medskip
Since the generating set is periodic with period $M$, the sequence of differences $a_{n+1} - a_n$ is purely periodic. Let $T = |\mathcal{S}^*|$ be the number of valid residue classes in one period of length $M$. It immediately follows that traversing $T$ terms in the sequence corresponds precisely to advancing by one full period $M$.

\medskip
Therefore, for all $n \geqslant N$, we have:
$$ \boxed{a_{n+T} = a_n + M} $$To ensure this holds for all positive integers $n$ as stipulated, observe that the sequence is strictly determined from the beginning. The aforementioned stabilisation applies globally because the elements of $\mathcal{S}^*$ are those sharing primes with all sequence members, thereby demonstrating that $T$ and $L = M$ satisfy the relation uniformly for all $n\geqslant 1$. \hfill$\square$
\end{proof}
\newpage
\subsubsection{Solution 3}
Call a prime \textit{small} if it is at most $a_1$, and big otherwise. Let the small part and big part of a positive integer $n$ be the (unique) $a$ and $b$ such that $a$ has only small prime factors, $b$ has only big prime factors, and $ab=n$. Note that the given condition implies $\gcd(a_i,a_j)>1$ for all $i,j$.
\begin{claim}
    For all $i,j$, $\gcd(a_i,a_j)$ has a small prime factor.
\end{claim}
\begin{proof}
    Suppose otherwise, and let $j$ be minimal such that there exists $i\leq j$ with $\gcd(a_i,a_j)$ having small part $1$. Let $s$ and $r$ be the small and big parts of $a_i$, respectively. Then $\gcd(a_i,a_j)=\gcd(r,a_j)$ is divisible by at least one large prime, so $r>a_1$, and also $\gcd(s,a_j)=1$. Furthermore, for all $k<j$, since $\gcd(a_i,a_k)$ has a small prime factor, we need $\gcd(s,a_k)>1$.

\medskip
Now consider the sequence $s,s^2,s^3,\ldots$. There must exist $t$ such that $s^t \in (a_1,a_1s] \subset (a_1,sr)=(a_1,a_i)$. I claim that $s^t$ appears in $(a_i)$. Suppose otherwise; since $s^t>a_1$, there exists $i$ such that $a_p<s^t<a_{p+1}$. Clearly $p<i$, so $\gcd(s,a_i)>1$ for all $i \leq p$. Then $\gcd(s^t,a_i)>1$ for all $i \leq p$ as well, but then by definition $a_{p+1}$ should be at most $s^t$: contradiction.

\medskip
Thus $s^t$ appears in $(a_i)$, and its index is evidently less than $i$. Therefore we would need $\gcd(s^t,a_j)>1$, but from before we have $\gcd(s,a_j)=1$: contradiction. Hence the claim is proved. \hfill$\square$
\end{proof}
\noindent
Now let $P$ be the product of all the small primes.
\begin{claim}
    For all $x,y \geq a_1$ with $P \mid x-y$, we have $x \in (a_i) \implies y \in (a_i)$
\end{claim}
\begin{proof}
    If $x \in (a_i)$, then from claim 1 $x$ and $a_i$ share a small prime factor for all $i$. Suppose $y \not \in (a_i)$, so there exists $p$ such that $a_p<y<a_{p+1}$. Because the small prime factors of $x$ and $y$ are the same, $y$ and $a_i$ also share a small prime factor for all $i$, implying $\gcd(y,a_i)>1$ for all $i \leq p$. But then $a_{p+1} \leq y$ by definition, which is a contradiction. Therefore $y \in (a_i)$. \hfill$\square$
\end{proof}
\noindent
\textcolor{GlacierBlue}{\textbf{Claim 2.3.8}} implies that there exists $T \geq 1$ with $a_{1+T}=a_1+P$, and then it is straightforward to show inductively that $a_{n+T}=a_n+P$ for all $n$, which is our desired result. $\square$
\begin{remark*}
To summarize the main/motiating idea of this solution: if we have some leading terms $a_1,\ldots,a_N$ where the gcd between any two elements is small (for some definition of small tbd), then if $s$ is the small part of $a_i$, every appropriately-sized multiple of $\operatorname{rad}(s)$ should also appear in $a_1,\ldots,a_N$. The rest is just filling in the details.
\end{remark*}
\begin{remark*}
A ``baby" version of the above remark is the following fact, which is easily seen to be true: if $x$ appears in $(a_i)$, then every multiple of $x$ also appears in $(a_i)$.
\end{remark*}
\begin{remark*}
Aside from $a_1=15$, I found the example of $a_1=75$ to be helpful. In this case we have
$$(a_i)=(75,78,80,84,90,96\ldots).$$Intuitively, it would be kind of surprising if there was some $j$ where the smallest prime dividing $\gcd(78,a_j)$ was $13$. This turns out to be easy to prove: if so, then $\gcd(6,a_j)=1$, but then $\gcd(96,a_j)=1$ as well. The problem here seems to stem from the fact that $13$ is a "big" prime, seeing as the prime divisors of $a_1$ are all at most $5$, which leads us down the right path.
\end{remark*}
\newpage
\subsubsection{Solution 4}
Let $C$ be the product of all primes not exceeding $2a_1$.
\begin{claim}
    For all $i, j$, we have $\gcd(a_i, a_j, C) > 1$.
\end{claim}
\begin{proof}
    Suppose on the contrary, and pick $(i, j)$ to be the smallest pair. Suppose there is some prime $p$ such that $p \nmid C$ and $p \mid a_i, a_j$. Hence $p\gcd(a_i, C) \mid a_i$. Since $p > a_i$, we have
\[
\gcd\left(p\gcd(a_i, C), a_k\right) > 1,\qquad \forall k = 1, 2, \ldots, i-1,
\]we must have $p\gcd(a_i, C)$ belonging to the sequence $(a_n)$. But
\[
\gcd\left(p\gcd(a_i, C), a_j, C\right) = 1,
\]hence $p\gcd(a_i, C) = a_i$. Since
\[
\dfrac{a_i}{a_1} > \gcd(a_i, C),
\]there exists some $k$ such that $\gcd(a_i, C)^k$ lies between $a_1$ and $a_i$, hence $\gcd(a_i, C)^k$ belongs to the sequence. But $\gcd(a_i, C)^k$ cannot be $a_i$, and
\[
\gcd\left(\gcd(a_i, C)^k, a_j, C\right) = 1,
\]a contradiction. \hfill$\square$
\end{proof}
\noindent
Let
\[
S=\left\{\gcd(a_i,C)\mid i\in\mathbb{Z}_{>0}\right\},
\]if there is some positive integer $m$ such that $\gcd(m,C)\in S$, then we must have
\[
\gcd(m,a_i)>1,\qquad \forall i\in\mathbb{Z}_{>0}.
\]Hence $m\in(a_n)$.

Hence this sequence is just the sequence of all numbers modulo $C$, therefore there exist $T$ and $L$ such that
\[
a_{n+T}=a_n+L.
\]
\newpage
\subsubsection{Solution 5}
For all $n$, let $A_n$ denote the set of all minimal (for the divisibility ordering) squarefree positive integers $x$ such that
$\gcd(a_i,x)>1$ for all $i\in [\![1,n]\!]$. This gives the following characterization :
\[\gcd(a_i,t)>1 \quad \text{for all $i\in [\![1,n]\!]$}\iff \exists x\in A_n, x\mid t.\]Furthermore, for a set $A\subset \mathbb N$, let $p(A)$ denote the set of all primes dividing at least one element of $A$. Clearly, each $p(A_n)$ is finite as $p(A_n)\subseteq p(\{a_1,a_2,\ldots,a_n\})$ and thus each $A_n$ is also finite. Let $N$ be an integer such that for $n>N$, $a_n>a_1^2$.
\begin{claim}
    For all $n>N$, we have $p(A_{n+1})\subseteq p(A_n)$.
\end{claim}
\begin{proof}
    Let $n>N$ be an integer. Let $x\in A_n$ be an element of $A_n$ which is at most $a_n$.

    \medskip
If $x\ge a_1$, then by definition $x$ must appear somewhere in the sequence $(a_n)$ between $a_1$ and $a_n$. Therefore, every element of $A_n$ has a non trivial $\gcd$ with $x$.
Otherwise, we have $x<a_1$. In that case, since $a_n>a_1^2$, there must be some power of $x$ between $a_1$ and $a_n$. By the same argument, this power must appear in the sequence and thus every element of $A_n$ has a non trivial $\gcd$ with $x$ in both cases.

\medskip
By definition, $a_{n+1}$ must be divisible by some element of $A_n$. If $x\mid a_{n+1}$ for some $x\in A_n\cap [\![1,a_n]\!]$, then
$A_{n+1}=A_n$ by our result above.
Otherwise, suppose there is some $y\in A_n$ with $y>a_n$ such that $y\mid a_{n+1}$. Since each multiple of $a_1$ appears in the sequence, we must have
\[a_{n+1}\le a_n+a_1\le 2a_n<2y,\]so that $a_{n+1}=y$. Then, since $p(\{a_{n+1}\})\subseteq p(A_n)$, we have $p(A_{n+1})\subseteq p(A_n)$. This concludes the proof of the lemma. \hfill$\square$
\end{proof}
\noindent
The lemma shows that the set $P$ of primes dividing some element of $A_n$ for some $n$ is finite. Let
\[M=\prod_{p\in P}p.\]Let $i\le j$ be positive integers. By definition, $a_j$ must be divisible by some element of $A_i$. That element of $A_i$ has a non trivial $\gcd$ with $a_i$ consisting of primes dividing $M$. In particular, $\gcd(a_i,a_j)$ has a prime factor dividing $M$.

The idea to finish is to consider the sequence modulo $M$ : let $0\le b_1<b_2<\ldots<b_k\le M-1$ be the remainders of elements of the sequence $(a_n)$ modulo $M$. Any two positive integers congruent to $b_i$ and $b_j$ modulo $M$ have a non trivial $\gcd$ (they share a prime factor dividing $M$). Therefore, we have the following characterization of the sequence $(a_n)$ :
\[a_{n+1}=\begin{cases}
        a_n+b_{i+1}-b_i \quad \text{if $a_n\equiv b_i\pmod M$ with $1\le i<k$}\\
        a_n+M-b_k+b_1\quad \text{if $a_n\equiv b_k\pmod M$}.
    \end{cases}\]We can thus conclude by taking $T=k$ and $L=M$. \hfill$\square$
\newpage
\subsubsection{Solution 6}
Let $$S = \{a_i \mid i \in \mathbb{N}\}.$$
Define $$b_{i,j} = \min_{n \in \mathbb{N}, n \mid \gcd(a_i,a_j)} n.$$
Let $s_x$ be the product of all prime not exceeding $$\max_{i,j \leqslant x} b_{i,j}.$$
\begin{claim}
    Let $n \in \mathbb{N}$. If $x \in \mathbb{N}$ satisfy $a_1 \leqslant x \leqslant a_n$ and $\gcd(x,a_i) >1$ for $i=1,2,\dots,n$, then $x \in S$.
\end{claim}
\begin{proof}
    Suppose that $x \not\in S$, this mean there exist $i$ such that $a_{i-1} < x < a_i$. But $\gcd(x,a_k) >1$ for $k=1,2,\dots,i-1$. Hence, $a_i$ is not the smallest positive integer greater than $a_{i-1}$ such that $\text{gcd}(a_{i}, a_k)>1$ for every $k=1,2,\ldots, i-1$. a contradiction.\hfill$\square$
\end{proof}
\begin{claim}
    Let $n<x$ be a positive integers. If $\gcd(a_n,s_x) =d$, then any multiple of $d$ between $a_1,a_x$ is an element of $S$.
\end{claim}
\begin{proof}
    Suppose that for some $i \leqslant x$, we have $\gcd(a_i,d) =1$, then $a_i$ and $d$ does not share any prime factor. But then $a_n$ contain all prime factor that $d$ has, so $\gcd(a_n,a_i)$ must only contain a prime factor $P$ that is coprime to $d$. But we know that $\text{rad}(P) \mid s_x$, so we get a contradiction. This mean $\gcd(d,a_i) >1$ for all $i \leqslant x$. Thus, by the claim above, we get this claim.\hfill$\square$
\end{proof}
\begin{claim}
    $b_{i,j} \leqslant a_1$.
\end{claim}
\begin{proof}
    Let $n$ be the smallest integer such that for some $m<n$, we have $b_{m,n} = P > a_1$
Note that by minimality of $n$,$P \nmid s_{n-1}$. Let $\gcd(a_i,s_{n-1}) = q_i$. Since $b_{m,n} =P$, we get that $q_m,a_n$ are coprime and so does $q_m,a_n$. Note that if $a_1$ is of the form $q_m^k$, then $\gcd(a_1,a_n) =1$, a contradiction. Hence suppose that $q_m^{k} > a_1 > q_m^{k-1}$. Note that$$P \mid \frac{a_m}{q_m} \implies P \leqslant \frac{a_m}{q_m} \implies q_m^{k-1} < a_1 < P \leqslant \frac{a_m}{q_m}$$Thus, we get that $a_1< q_m^k = q_m \cdot q_m^{k-1} < q_m \cdot \frac{a_m}{q_m} = a_m$. Thus, $q_m^k \in S$. However, $\gcd(a_n,q_m^k) = 1$, a contradiction.\hfill$\square$
\end{proof}
\noindent
Let $Q$ be the product of primes not exceeding $a_1$. Let $R = \{x \mid x \in S, x < Q+a_1\}$.
\begin{claim}
    $k \in S \iff Q \mid k-i$ for some $i \in R$.
\end{claim}
\begin{proof}
    Suppose $k \in S$, then by claim $2$, all multiple of $\gcd(k,Q)$ which is at least $a_1$ is an element of $S$. This mean $k \pm Q \in S$. This clearly finish the claim.\hfill$\square$
\end{proof}
\noindent
Hence, we get that $a_{n+|R|} = a_{n} + Q$ as desired.

\end{document}
